\chapter{A comparison geometry approach to the Ricci flow}\label{chap:comparison-geometry}
\label{lengthfn}

For an $n$-dimensional generalized Ricci flow, this chapter develops
Perelman's ${\mathcal L}$-geometry: critical curves, the
${\mathcal L}$-exponential map, reduced length, and reduced volume
\cite{Perelman1}.


\section{${\mathcal L}$-length and ${\mathcal L}$-geodesics}

Fix an $n$-dimensional generalized Ricci flow $({\mathcal M},G)$ and
a time $T$ distinct from its initial time. Write
$T{\mathcal M}=\langle\chi\rangle\oplus {\mathcal HT}{\mathcal M}$.


\begin{definition}[Parameterized by backward time]\label{def:parameterized-by-backward-time}
\uses{def:generalized-ricci-flow, def:time-slices-horizontal}

Let $0\le \tau_1<\tau_2$ be given and let $\gamma\colon
[\tau_1,\tau_2]\to {\mathcal M}$ be a continuous map. We say that
$\gamma$ is {\em parameterized by backward time} provided that $\gamma(\tau)\in M_{T-\tau}$ for all
$\tau\in [\tau_1,\tau_2]$.
\end{definition}


Throughout this section the paths $\gamma$ that we consider shall be
parameterized by backward time. We begin with the definition of
${\mathcal L}$-length of such a path.

\begin{definition}[${\mathcal L}$-length]\label{def:l-length}
\uses{def:parameterized-by-backward-time, def:ricci-curvature}
\label{Llength}

Let  $\gamma\colon [\tau_{1},\tau_{2}]\rightarrow {\mathcal M},\ 0
\le \tau_{1} < \tau_{2}$, be a  $C^{1}$-path parameterized by
backward time. We define $X_\gamma(\tau)$ to be the horizontal
projection of the tangent vector $d\gamma(\tau)/d\tau$, so that
$d\gamma/d\tau=-\chi+X_\gamma(\tau)$ with $X_\gamma(\tau)\in
{\mathcal HT}{\mathcal M}$. We define the $\mathcal{L}$-{\em
length} of $\gamma$ to be:
$$\mathcal{L}(\gamma) = \int_{\tau_{1}}^{\tau_{2}}
\sqrt{\tau}\left(R(\gamma(\tau)) +
\abs{X_\gamma(\tau)}^{2}\right)d\tau,$$
where the norm of $X_\gamma(\tau)$ is measured using the metric
$G_{T-\tau}$ on ${\mathcal HT}{\mathcal M}$. When $\gamma$ is clear
from the context, we write $X$ for $X_\gamma$.
\end{definition}



With a view toward better understanding the properties of the paths
that are critical points of this functional, the so-called
$\mathcal{L}$-geodesics, especially near $\tau=0$, it is helpful to
introduce a convenient reparameterization. We set $s=\sqrt{\tau}$.
We use the notation $A(s)$ to denote the horizontal component of the
derivative of $\gamma$ with respect to the variable $s$. One sees
immediately by the chain rule that
\begin{equation}\label{changeofvar}
A(s^2) = 2s X(s^2)\ \ \text{or} \ \
A(\tau)=2\sqrt{\tau}X(\tau).\end{equation} With respect to the
variable $s$, the ${\mathcal L}$-functional is
\begin{equation}\label{seqn}
\mathcal{L}(\gamma)=
\int_{\sqrt{\tau_1}}^{\sqrt{\tau_2}}\left(\frac{1}{2}|A(s)|^{2} +
2R(\gamma(s))s^{2}\right)ds.
\end{equation}

Let's consider the simplest example.

\begin{example}[${\mathcal L}$-length in flat Euclidean space]\label{ex:l-length-euclidean}
\uses{def:l-length}

Suppose that our generalized Ricci flow is a constant family of
Euclidean metrics on $\Ar^n\times [0,T]$. That is to say, $g(t) =
g_{0}$ is the usual Euclidean metric. Then we have
$R(\gamma(\tau))\equiv 0$. Using the change of variables
$s=\sqrt{\tau}$, we have
$$\mathcal{L}(\gamma) = \frac{1}{2}\int_{\sqrt{\tau_1}}^{\sqrt{\tau_2}} \abs{A(s)}^{2}ds,
$$ which is the standard energy functional in Riemannian geometry
for the path $\gamma(s)$. The minimizers for this functional are the
maps $s\mapsto (\alpha(s),T-s^2)$ where $\alpha(s)$ is a straight
line in $\Ar^n$ parameterized at constant speed. Written in the
$\tau$ variables the minimizers are
$$\gamma(\tau)=(x+\sqrt{\tau}v,T-\tau),$$
straight lines parameterized at speed varying linearly with
$\sqrt{\tau}$.

\end{example}

\subsection{${\mathcal L}$-geodesics}

\begin{lemma}[Euler-Lagrange equation for ${\mathcal L}$-length]\label{lem:l-length-euler-lagrange}
\uses{def:l-length}
\label{EL}

The
Euler-Lagrange equation for critical paths for the $\mathcal{L}$-length is
\begin{equation}\label{EulerLag}
 \nabla_{X}X - \frac{1}{2}\nabla R
+ \frac{1}{2\tau}X + 2\Ric(X,\cdot)^* = 0.
\end{equation}
\end{lemma}


\begin{remark}[$\Ric(X,\cdot)^*$ as a horizontal vector field]\label{rem:ricci-one-form-dual}
\uses{lem:l-length-euler-lagrange}

$\Ric(X,\cdot)$ is a horizontal one-form along $\gamma$ and its
dual $\Ric(X,\cdot)^*$ is a horizontal tangent vector field
along $\gamma$.
\end{remark}

\begin{proof}
\sketch
\uses{def:l-length,def:ricci-flow}
First, let us suppose that the
generalized Ricci flow is an ordinary Ricci flow $(M,g(t))$. Let
$\gamma_u(\tau)=\gamma(\tau,u)$ be a family of curves parameterized
by backward time. Let
$$\widetilde Y(\tau,u)=\frac{\partial\gamma}{\partial u}.$$ Then $\widetilde X(\tau,u)=X_{\gamma_u}(\tau,u)$
 and $\widetilde Y(\tau,u)$
are  the coordinate vector fields along the surface obtained by taking the
projection of $\gamma(\tau,u)$ into $M$. Thus, $[\widetilde X,\widetilde Y]
=0$. We denote by $X$ and $Y$ the restrictions of $\widetilde X$ and
$\widetilde Y$, respectively to $\gamma_0$. We have
\begin{align*}
\frac{d}{du}\mathcal{L}(\gamma_{u})\bigl|_{u=0}\bigr.  & =
\frac{d}{du}\left(\int_{\tau_{1}}^{\tau_{2}}
\sqrt{\tau}(R(\gamma_u(\tau)) + \abs{\widetilde
X(\tau,u)}^{2})d\tau\right)\Bigl|_{u=0}\Bigr.
\\&=\int_{\tau_{1}}^{\tau_{2}} \sqrt{\tau}(\langle \nabla R, Y\rangle  +
2\langle (\nabla_{Y}\widetilde X)|_{u=0},X\rangle )d\tau
\end{align*}
On the other hand, since $\partial g/\partial \tau =2\Ric$ and
since $[\widetilde X,\widetilde Y]=0$, we have
\begin{eqnarray*} 2\frac{d}{d\tau}(\sqrt{\tau}\langle Y,X\rangle_{g(T-\tau)} )
& = & \frac{1}{\sqrt{\tau}}\langle Y,X\rangle +
2\sqrt{\tau}\langle\nabla_XY,X\rangle
+2\sqrt{\tau}\langle Y,\nabla_XX\rangle \\ & & +4\sqrt{\tau}\Ric(Y,X) \\
& = & \frac{1}{\sqrt{\tau}}\langle Y,X\rangle  +
2\sqrt{\tau}\langle(\nabla_Y\widetilde X)|_{u=0},X\rangle
+2\sqrt{\tau}\langle Y,\nabla_XX\rangle \\ & &  +4\sqrt{\tau}\Ric(Y,X)
\end{eqnarray*}
Using this we obtain \begin{eqnarray}\nonumber \frac{d}{du}
\mathcal{L}(\gamma_u)\bigl|_{u=0}\bigr. & = & \int_{\tau_1}^{\tau_2}
\Bigl(2\frac{d}{d\tau}\left[(\sqrt{\tau})\langle
Y,X\rangle\right]-\frac{1}{\sqrt{\tau}}\langle
Y,X\rangle\Bigr. \\
& & +\sqrt{\tau}\bigl(\langle\nabla R,Y\rangle
 - 2\langle Y,\nabla_{X}X\rangle  - 4\Ric(X,Y)\bigr)\Bigr)d\tau \nonumber \\
 & = & 2\sqrt{\tau}\langle Y,X\rangle |_{\tau_{1}}^{\tau_{2}} \nonumber \\
& & + \int_{\tau_{1}}^{\tau_{2}} \sqrt{\tau}\langle Y,\bigl(\nabla R
- \frac{1}{\tau}X - 2\nabla_{X}X - 4\Ric(X,\cdot)^*\bigr)\rangle d\tau.\label{variation}
\end{eqnarray}

Now we drop the assumption that the generalized Ricci flow is an
ordinary Ricci flow. Still we can partition the interval
$[\tau_1,\tau_2]$ into finitely many sub-intervals with the property
that the restriction of $\gamma_0$ to each of the sub-intervals is
contained in a patch of space-time on which the generalized Ricci
flow is isomorphic to an ordinary Ricci flow. The above argument
then applies to each of the sub-intervals. Adding up
Equation~(\ref{variation}) over these sub-intervals shows that the
same equation for the first variation of length for the entire
family $\gamma_u$ holds.


We consider a variation $\gamma(\tau,u)$ with fixed endpoints, so
that  $Y(\tau_{1})=Y(\tau_{2})= 0.$ Thus, the condition that
$\gamma$ be a critical path for the ${\mathcal L}$-length is that
the integral expression vanish for all variations $Y$ satisfying
$Y(\tau_1)=Y(\tau_2)=0$. Equation~(\ref{variation}) holds for all
such $Y$ if and only if $\gamma$ satisfies
Equation~(\ref{EulerLag}).
\end{proof}

\begin{remark}[Horizontal gradient convention]\label{rem:horizontal-gradient-convention}
\uses{lem:l-length-euler-lagrange}

In the Euler-Lagrange equation, $\nabla R$ is the horizontal
gradient, and the equation is an equation of horizontal vector
fields along $\gamma$.
\end{remark}

\begin{definition}[${\mathcal L}$-geodesic]\label{def:l-geodesic}
\uses{def:l-length, lem:l-length-euler-lagrange}

A curve $\gamma$, parameterized by backward time, that is a critical
point of the ${\mathcal L}$-length is called an $\mathcal{L}$-{\em
geodesic}.
Equation~(\ref{EulerLag}) is the {\em ${\mathcal L}$-geodesic
equation}.
\end{definition}


Written with respect to the variable $s=\sqrt{\tau}$ the
$\mathcal{L}$-geodesic equation becomes
\begin{equation}\label{reparam}
\nabla_{A(s)}A(s) -2s^{2}\nabla R + 4s\Ric(A(s),\cdot)^*=0.\end{equation}
 Notice that in this form the ODE is
regular even at $s=0$.



\begin{lemma}[Initial limit determines an ${\mathcal L}$-geodesic]\label{lem:l-geodesic-initial-limit}
\uses{def:l-geodesic}

Let $\gamma\colon [0,\tau_2]\to {\mathcal M}$ be an ${\mathcal L}$-geodesic.
Then $\lim_{\tau\rightarrow 0}\sqrt{\tau}X_\gamma(\tau)$ exists. The
${\mathcal L}$-geodesic $\gamma$ is completely determined by this limit (and by
$\tau_2$).
\end{lemma}

\begin{proof}
\sketch
\uses{def:l-geodesic}
Since the ODE in Equation~(\ref{reparam})  is non-singular even at
zero, it follows that $A(s)$ is a smooth function of $s$ in a
neighborhood of $s=0$. , The lemma follows easily by the change of
variables formula, $A(\tau)=2\sqrt{\tau}X_\gamma(\tau)$.
\end{proof}

\begin{definition}[Minimizing ${\mathcal L}$-geodesic]\label{def:l-geodesic-minimizing}
\uses{def:l-geodesic, def:l-length}

An ${\mathcal L}$-geodesic is said to be {\em
minimizing} if there is no
curve parameterized by backward time with the same endpoints and
with smaller ${\mathcal L}$-length.
\end{definition}


\subsection{The ${\mathcal L}$-Jacobi equation}

Consider a family $\gamma(\tau,u)$ of ${\mathcal L}$-geodesics parameterized by
$u$ and defined on $[\tau_1,\tau_2]$ with $0\le \tau_1<\tau_2$. Let $Y(\tau)$
be the horizontal vector field along $\gamma$ defined by
$$Y(\tau)=\frac{\partial}{\partial u}\gamma(\tau,u)|_{u=0}.$$


\begin{lemma}[The ${\mathcal L}$-Jacobi equation]\label{lem:l-jacobi-equation}
\uses{def:l-geodesic, lem:l-length-euler-lagrange}

$Y(\tau)$ satisfies the {\em ${\mathcal L}$-Jacobi
equation}:
\begin{equation}\label{jaceqn}\nabla_X\nabla_XY+{\mathcal
R}(Y,X)X-\frac{1}{2}\nabla_Y(\nabla
R)+\frac{1}{2\tau}\nabla_XY+2(\nabla_Y\Ric)(X,\cdot)^*+2\Ric(\nabla_XY,\cdot)^*=0.\end{equation} This is a second-order
linear equation for $Y$. Supposing that $\tau_1>0$, there is a
unique horizontal vector field $Y$ along $\gamma$ solving this
equation vanishing at $\tau_1$ with a given first-order derivative
along $\gamma$ at $\tau_1$. Similarly, there is a unique solution
$Y$ to this equation vanishing at $\tau_2$ and with a given
first-order derivative at $\tau_2$.
\end{lemma}


\begin{proof}
\sketch
\uses{lem:l-length-euler-lagrange}
Given a family $\gamma(\tau,u)$ of ${\mathcal L}$-geodesics, then
from Lemma~\ref{lem:l-length-euler-lagrange} we have
$$ \nabla_{\widetilde X}\widetilde X = \frac{1}{2}\nabla R(\gamma)
-\frac{1}{2\tau}\widetilde X - 2\Ric(\widetilde X,\cdot)^*.$$
Differentiating this equation in the $u$-direction along the curve
$u=0$ yields
$$\nabla_Y
\nabla_{\widetilde X}\widetilde X|_{u=0}=\frac{1}{2}\nabla_Y(\nabla
R)-\frac{1}{2\tau}\nabla_Y(\widetilde X)|_{u=0} -2\nabla_Y(\Ric(\widetilde X,\cdot))^*|_{u=0}.$$ Of course, we have
$$\nabla_Y(\Ric(\widetilde X,\cdot)^*)|_{u=0}=(\nabla_Y\Ric)(X,\cdot)^*
+\Ric(\nabla_Y\widetilde X|_{u=0},\cdot)^*.$$ Plugging this in,
interchanging the orders of differentiation on the left-hand side, using
$\nabla_{\widetilde Y}\widetilde X=\nabla_{\widetilde X}\widetilde Y$, and
restricting to $u=0$ yields the equation given in the statement of the lemma.
This equation is a regular, second-order linear equation  for all $\tau>0$, and
hence is determined by specifying the value and first derivative at any
$\tau>0$.
\end{proof}

Equation~(\ref{jaceqn}) is obtained by applying $\nabla_Y$ to
Equation~(\ref{EulerLag}) and exchanging orders of differentiation.
The result Equation~(\ref{jaceqn}) is a second-order differential
equation for $Y$ that makes no reference to an extension of
$\gamma(\tau)$ to an entire family of curves.



\begin{definition}[${\mathcal L}$-Jacobi field]\label{def:l-jacobi-field}
\uses{lem:l-jacobi-equation}

A field $Y(\tau)$ along an ${\mathcal L}$-geodesic is
called an {\em ${\mathcal L}$-Jacobi field} if it satisfies the ${\mathcal L}$-Jacobi
equation, Equation~(\ref{jaceqn}), and if it vanishes at $\tau_1$.
For any horizontal vector field $Y$ along $\gamma$ we denote by
$\Jac(Y)$ the expression on the left-hand
side of Equation~(\ref{jaceqn}).
\end{definition}

In fact, there is a similar result even for $\tau_1=0$.

\begin{lemma}[Initial limit determines an ${\mathcal L}$-Jacobi field]\label{lem:l-jacobi-field-initial-limit}
\uses{def:l-jacobi-field}

Let $\gamma$ be an ${\mathcal L}$-geodesic defined on $[0,\tau_2]$
and let $Y(\tau)$ be an ${\mathcal L}$-Jacobi field along $\gamma$.
Then
$$\lim_{\tau\rightarrow 0}\sqrt{\tau}\nabla_XY$$
exists. Furthermore, $Y(\tau)$ is completely determined by this
limit.
\end{lemma}


\begin{proof}
\sketch
\uses{def:l-jacobi-field}
We use the variable $s=\sqrt{\tau}$, and let $A(s)$ be the
horizontal component of $d\gamma/ds$. Then differentiating the
${\mathcal L}$-geodesic equation written with respect to this
variable we see
$$\nabla_A\nabla_AY=-{\mathcal R}(Y,A)A+2s^2\nabla_Y(\nabla R)-4s(\nabla_Y\Ric)(A,\cdot)
-4s\Ric(\nabla_AY,\cdot).$$ Hence, for each tangent vector $Z$,
there is a unique solution to this equation with the two initial
conditions $Y(0)=0$ and $\nabla_AY(0)=Z$.

On the other hand, from Equation~(\ref{changeofvar}) we have
$\nabla_X(Y)=\frac{1}{2\sqrt{\tau}}\nabla_A(Y)$, so that
$$\sqrt{\tau}\nabla_X(Y)=\frac{1}{2}\nabla_A(Y).$$
\end{proof}


\subsection{Second order variation of ${\mathcal
L}$} We shall need the relationship of the ${\mathcal L}$-Jacobi
equation to the second-order variation of ${\mathcal L}$. This is
given in the next proposition.




\begin{proposition}[Second variation of ${\mathcal L}$-length via ${\mathcal L}$-Jacobi fields]\label{prop:l-length-second-variation}
\uses{def:l-geodesic-minimizing, def:l-jacobi-field}
\label{2ndvarj}

Suppose that $\gamma$ is a minimizing ${\mathcal L}$-geodesic. Then,
for any vector field $Y$ along $\gamma$, vanishing at both
endpoints, and any family $\gamma_u$ of curves parameterized by
backward time with $\gamma_0=\gamma$ and with the $u$-derivative of
the family at $u=0$ being the vector field $Y$ along $\gamma$, we
have
$$\frac{d^2}{du^2}{\mathcal L}(\gamma_u)|_{u=0}=-\int_{\tau_1}^{\tau_2}2\sqrt{\tau}
\langle \Jac(Y),Y\rangle d\tau.$$ This quantity vanishes if and
only if $Y$ is an ${\mathcal L}$-Jacobi field.
\end{proposition}


Let us begin the proof of this proposition with the essential computation.


\begin{lemma}[Second variation cross-term formula]\label{lem:l-length-second-variation-formula}
\uses{def:l-jacobi-field, def:l-length}
\label{2ndvari}

Let $\gamma$ be an ${\mathcal L}$-geodesic defined on $
[\tau_1,\tau_2]$, and let $Y_1$ and $Y_2$ be horizontal vector
fields along $\gamma$ vanishing at $\tau_1$. Suppose that
$\gamma_{u_1,u_2}$ is any family of curves parameterized by backward
time with the property that $\gamma_{0,0}=\gamma$ and the derivative
of the family in the $u_i$-direction at $u_1=u_2=0$ is $Y_i$. Let
$\widetilde Y_i$ be the image of $\partial/\partial u_i$ under
$\gamma_{u_1,u_2}$ and let $\widetilde X$ be the image of the
horizontal projection of $\partial/\partial \tau$ under this same
map, so that the restrictions of these three vector fields to the
curve $\gamma_{0,0}=\gamma$ are $Y_1,Y_2$ and $X$ respectively. Then
we have \begin{eqnarray*} \frac{\partial}{\partial
u_1}\frac{\partial}{\partial u_2}{\mathcal
L}(\gamma_{u_1,u_2})|_{u_1=u_2=0} & = & 2\sqrt{\tau_2}
Y_1(\tau_2)\langle \widetilde Y_2(\tau_2,u_1,0),\widetilde
X(\tau_2,u_1,0)\rangle|_{u_1=0}\\
& & -\int_{\tau_1}^{\tau_2}2\sqrt{\tau}\langle \Jac(Y_1),Y_2\rangle d\tau.\end{eqnarray*}
\end{lemma}


\begin{proof}
\sketch
\uses{lem:l-length-euler-lagrange}
According to Equation~(\ref{variation}) we have
\begin{eqnarray*}
\frac{\partial}{\partial u_2}{\mathcal L}(\gamma)(u_1,u_2)
 & = & 2\sqrt{\tau_2}\langle \widetilde Y_2(\tau_2,u_1,u_2),\widetilde
X(\tau_2,u_1,u_2)\rangle
\\ & & -\int_{\tau_1}^{\tau_2}2\sqrt{\tau}\langle EL(\widetilde
X(\tau,u_1,u_2),\widetilde Y_2(\tau,u_1,u_2)\rangle d\tau,\end{eqnarray*} where
$EL(\widetilde X(\tau,u_1,u_2))$ is the Euler-Lagrange expression for
geodesics, i.e., the left-hand side of Equation~(\ref{EulerLag}).
Differentiating again yields:
\begin{eqnarray}
\lefteqn{\frac{\partial }{\partial u_1}\frac{\partial }{\partial u_2}{\mathcal
L}(\gamma_{u_1,u_2}\bigl)|_{u_1=u_2=0}\big.
 =  2\sqrt{\tau_2}Y_1(\tau_2)\langle \widetilde
Y_2(\tau_2,u_1,0),\widetilde X(\tau_2,u_1,0)\rangle\bigl|_{u_1=0}\bigr.  } \nonumber\\
& & -
\int_{\tau_1}^{\tau_2}2\sqrt{\tau}\left(\langle\nabla_{Y_1}EL(\widetilde
X), Y_2\rangle +\langle EL(X),\nabla_{Y_1}\widetilde
Y_2\rangle\right)(\tau,0,0)d\tau. \label{crosspart} \end{eqnarray}
Since $\gamma_{0,0}=\gamma$ is a geodesic, the second term in the
integrand vanishes, and since $[\widetilde X,\widetilde Y_1]=0$, we
have $\nabla_{Y_1}EL(\widetilde X(\tau,0,0))=\Jac(Y_1)(\tau)$.
This proves the lemma.
\end{proof}

\begin{remark}[Second variation is determined by the ${\mathcal L}$-Jacobi field]\label{rem:second-variation-determined-by-jacobi}
\uses{lem:l-length-second-variation-formula, def:l-jacobi-field}
\label{variremark}

Let $\gamma(\tau,u)$ be a family of curves as above with
$\gamma(\tau,0),\ \tau_1\le \tau\le \bar \tau$, being an ${\mathcal
L}$-geodesic. It follows from Lemma~\ref{lem:l-length-second-variation-formula} and the remark
after the introduction of the ${\mathcal L}$-Jacobi equation that
the second-order variation of length at $u=0$ of this family is
determined by the vector field $Y(\tau)=\partial \gamma/\partial u$
along $\gamma(\cdot,0)$ and by the second-order information about
the curve $\gamma(\bar\tau,u)$ at $u=0$.
\end{remark}


\begin{corollary}[Symmetry of the ${\mathcal L}$-Jacobi bilinear pairing]\label{cor:l-jacobi-bilinear-symmetric}
\uses{def:l-jacobi-field, def:l-geodesic}
\label{symmY}

Let $\gamma$ be an ${\mathcal L}$-geodesic and let $Y_1,Y_2$ be vector fields
along $\gamma$ vanishing at $\tau_1$. Suppose $Y_1(\tau_2)=Y_2(\tau_2)=0$.
 Then the
bilinear pairing
$$-\int_{\tau_1}^{\tau_2}2\sqrt{\tau}\langle
\Jac(Y_1),Y_2\rangle d\tau$$ is a symmetric function of $Y_1$
and $Y_2$.
\end{corollary}

\begin{proof}
  \uses{lem:l-length-second-variation-formula, cor:l-jacobi-bilinear-symmetric, def:l-geodesic-minimizing}
\sketch

Given $Y_1$ and $Y_2$ along $\gamma$ we construct a two-parameter
family of curves parameterized by backward time as follows. Let
$\gamma(\tau,u_1)$ be the value at $u_1$ of the geodesic through
$\gamma(\tau)$ with tangent vector $Y_1(\tau)$. This defines a
family of curves parameterized by backward time, the family being
parameterized by  $u_1$ sufficiently close to $0$. We extend
 $Y_1$ and $X$ to vector fields on this entire family by defining them to be
 $\partial/\partial u_1$ and the
horizontal projection of $\partial/\partial \tau$, respectively.
 Now we
extend the vector field $Y_2$ along $\gamma$ to a vector field on
this entire one-parameter family of curves. We do this so that
$Y_2(\tau_2,u_1)=Y_1(\tau_2,u_1)$. Now given this extension
$Y_2(\tau,u_1)$ we define a two-parameter family of curves
parameterized by backward time
 by setting $\gamma(\tau,u_1,u_2)$ equal to the value at $u_2$ of the geodesic
 through $\gamma(\tau,u_1)$ in the direction $Y_2(\tau,u_1)$.
 We then extend  $Y_1$, $Y_2$, and $X$ over this entire family by letting them be
 $\partial/\partial u_1$,
$\partial/\partial u_2$, and the horizontal projection of $\partial/ \partial
\tau$, respectively. Applying Lemma~\ref{lem:l-length-second-variation-formula} and using the fact that
$Y_i(\bar\tau)=0$ we conclude that
$$\frac{\partial}{\partial u_1}\frac{\partial}{\partial u_2}{\mathcal L}(\gamma)|_{u_1=u_2=0}= -\int_{\tau_1}^{\tau_2}2\sqrt{\tau}\langle
\Jac(Y_1),Y_2\rangle d\tau$$ and symmetrically that
$$\frac{\partial}{\partial u_2}\frac{\partial}{\partial u_1}{\mathcal L}(\gamma)|_{u_1=u_2=0}= -\int_{\tau_1}^{\tau_2}2\sqrt{\tau}\langle
\Jac(Y_2),Y_1\rangle d\tau.$$ Since the second cross partials
are equal, the corollary follows.
\end{proof}





Now we are in a position to establish Proposition~\ref{prop:l-length-second-variation}.


\begin{proof}
\sketch
\uses{lem:l-length-second-variation-formula,cor:l-jacobi-bilinear-symmetric,def:l-geodesic-minimizing}
From the equation in Lemma~\ref{lem:l-length-second-variation-formula}, the equality of the second variation
of ${\mathcal L}$-length at $u=0$ and the integral is immediate from the fact
that $Y(\tau_2)=0$. It follows immediately that, if  $Y$ is an ${\mathcal
L}$-Jacobi field vanishing at $\tau_2$, then the second variation of the length
vanishes at $u=0$. Conversely, suppose given a family $\gamma_u$ with
$\gamma_0=\gamma$ with the property that the second variation of length
vanishes at $u=0$, and that the vector field $Y=(\partial\gamma/\partial
u)|_{u=0}$ along $\gamma$ vanishes at the end points. It follows that the
integral also vanishes. Since $\gamma$ is a minimizing ${\mathcal L}$-geodesic,
for any variation $W$, vanishing at the endpoints, the first variation of the
length vanishes and the second variation of length is non-negative. That is to
say,
$$-\int_{\tau_1}^{\tau_2}2\sqrt{\tau}\langle \Jac(W),W\rangle d\tau\ge 0$$
for all vector fields $W$ along $\gamma$ vanishing at the endpoints.
Hence, the restriction to the space of vector fields along $\gamma$
vanishing at the endpoints of the bilinear form
$$B(Y_1,Y_2)=-\int_{\tau_1}^{\tau_2}2\sqrt{\tau}\langle
\Jac(Y_1),(Y_2)d\tau,$$ which is symmetric by Corollary~\ref{cor:l-jacobi-bilinear-symmetric}, is
positive semi-definite. Since $B(Y,Y)=0$, it follows that $B(Y,\cdot)=0$; that
is to say, $\Jac(Y)=0$.
\end{proof}


\section{The ${\mathcal L}$-exponential map and its first-order
properties}\label{sect:Lexp}

We use  ${\mathcal L}$-geodesics in order to define the ${\mathcal
L}$-exponential map.

For Section~\ref{sect:Lexp} we fix $\tau_1\ge 0$ and a point $x\in
{\mathcal M}$ with ${\bf t}(x)=T-\tau_1$. We suppose that $T-\tau_1$
is greater than the initial time of the generalized Ricci flow.
Then, for every $Z \in T_{x}M_{T-\tau_1},$ there is a maximal
$\mathcal{L}$-geodesic, denoted $\gamma_{Z}$, defined on some
positive $\tau$-interval, with $\gamma_{Z}(\tau_1)=x$ and with
$\sqrt{\tau_1}X(\tau_1)=Z$. (In the case $\tau_1=0$, this equation
is interpreted to mean $\lim_{\tau\rightarrow 0} \sqrt{\tau}X(\tau)
= Z$.)

\begin{definition}[${\mathcal L}$-exponential map]\label{def:l-exponential-map}
\uses{def:l-geodesic}
\label{defnD}

We define {\em the domain of definition of} ${\mathcal L}\mathrm{exp}_x$,  denoted ${\mathcal D}_x$, to be the
subset of $T_xM_{T-\tau_1}\times (\tau_1,\infty)$ consisting of all pairs
$(Z,\tau)$ for which $\tau> \tau_1$ is in the maximal domain of definition of
$\gamma_Z$.  Then we define $\mathcal{L}\exp_{x}\colon {\mathcal D}_x\to
{\mathcal M}$ by setting $\mathcal{L}\exp_{x}(Z,\tau) = \gamma_{Z}(\tau)$
for all $(Z,\tau)\in {\mathcal D}_x$. We
define the map $\widetilde L\colon {\mathcal D}_x\to \Ar$
by $\widetilde L(Z,\tau)={\mathcal L}\left(\gamma_Z|_{[\tau_1,\tau]}\right)$.
Lastly, for any $\tau> \tau_1$ we denote by ${\mathcal L}\mathrm{exp}_x^\tau$ the restriction of
${\mathcal L}\exp_x$ to the slice
$${\mathcal
D}^\tau_x={\mathcal D}_x\cap \left(T_xM_{T-\tau_1}\times \{\tau\}\right),$$
which is {\em the domain of definition of} ${\mathcal L}\exp_x^\tau$. We
also denote by  $\widetilde L^\tau$ the
restriction of $\widetilde L$ to this slice. We will implicitly identify
${\mathcal D}_x^\tau$ with a subset of $T_xM_{T-\tau_1}$.
\end{definition}



\begin{lemma}[${\mathcal D}_x$ is open and ${\mathcal L}\exp_x$ is smooth]\label{lem:l-exponential-smooth}
\uses{def:l-exponential-map}
\label{lcdiffeo}

${\mathcal D}_x$ is an open subset of $T_xM_{T-\tau_1}\times (\tau_1,\infty)$;
its intersection with each line $\{Z\}\times (\tau_1,\infty)$ is a non-empty
interval whose closure contains $\tau_1$. Furthermore, ${\mathcal L}\mathrm{exp}_x\colon {\mathcal D}_x\to {\mathcal M}$ is a smooth map, and $\widetilde
L$ is a smooth function.
\end{lemma}

\begin{proof}
\sketch
\uses{lem:l-length-euler-lagrange}
The tangent vector in space-time of  the ${\mathcal L}$-geodesic
$\gamma$ is the vector field $-\chi+X_\gamma(\tau)$ along $\gamma$,
where $X_\gamma(\tau)$ satisfies Equation~(\ref{EulerLag}). As
above, in the case $\tau_1=0$, it is convenient to replace the
independent variable $\tau$ by $s=\sqrt{\tau}$, so that the ODE
becomes Equation~(\ref{reparam}) which is regular at $0$. With this
change, the lemma then follows immediately by the usual results on
existence, uniqueness and $C^\infty$-variation with parameters of
ODE's.
\end{proof}


\subsection{The differential of ${\mathcal L}\exp$}
Now we compute the differential of ${\mathcal L}\exp$.


\begin{lemma}[Differential of ${\mathcal L}\exp$ via ${\mathcal L}$-Jacobi fields]\label{lem:l-exponential-differential}
\uses{def:l-exponential-map, def:l-jacobi-field}
\label{DLJacobi}

Let $Z\in {\mathcal D}_x^{\bar\tau}\subset T_xM_{T-\tau_1}$. The
differential of ${\mathcal L}\mathrm{exp}^{\bar\tau}_x$
at the point $Z$ is given as follows: For each $W\in
T_x(M_{T-\tau_1})$ there is a unique ${\mathcal L}$-Jacobi
field $Y_W(\tau)$ along
$\gamma_Z$ with the property that $\sqrt{\tau_1}Y_W(\tau_1)=0$ and
$\sqrt{\tau_1}\nabla_X(Y_W)(\tau_1)=W$. We have
$$d_Z{\mathcal L}\exp^{\bar\tau}_x(W)=Y_W(\bar\tau).$$
Again, in case $\tau_1=0$, both of the conditions on $Y_W$ are interpreted as
the limits as $\tau\rightarrow 0$.
\end{lemma}


\begin{proof}
\sketch
\uses{def:l-jacobi-field}
Let $Z(u)$ be a curve in ${\mathcal D}_x^{\bar\tau}$ with $Z(0)=Z$.
Let $\gamma_u$ be the ${\mathcal L}$-geodesic starting at $x$ with
$\sqrt{\tau_1}X_{\gamma_u}(\tau_1)=Z(u)$. Then, clearly,
$$d_Z{\mathcal L}\exp^{\bar\tau}_x\left(\frac{dZ}{du}(0)\right)
=\frac{\partial}{\partial u}\left(\gamma_u(\bar\tau)\right)|_{u=0}.$$ On the
other hand, the vector field $Y(\tau)=\left(\partial\gamma_u(\tau)/\partial
u\right)|_{u=0} $ is an ${\mathcal L}$-Jacobi field along $\gamma_Z$.  Thus, to
complete the proof in the case when $\tau_1>0$ we need only see that
$\nabla_X\widetilde Y(\tau_1)=\nabla_Y\widetilde X(\tau_1)$. This is clear
since, as we have already seen, $[\widetilde X,\widetilde Y]=0$.

  When $\tau_1=0$, the substitution $s=\sqrt\tau$ gives the stated
limit identity.
\end{proof}

\begin{claim}[Tangential limit commutes with the Jacobi limit]\label{claim:l-exp-tau-derivative-limit}
If $\tau_1=0$, then
$$\frac{\partial}{\partial u}\left(\lim_{\tau\rightarrow 0}\sqrt{\tau}X(\tau,u)\right)|_{u=0}=
\lim_{\tau\rightarrow 0}\sqrt{\tau}\frac{d}{d\tau}Y(\tau).$$
\end{claim}

\begin{proof}
\sketch This follows by the substitution $s=\sqrt{\tau}$.
replacing $\tau$ by $s=\sqrt{\tau}$.
\end{proof}

\subsection{Positivity of the second variation at a minimizing
${\mathcal L}$-geodesic}


If $\gamma$ is a minimizing ${\mathcal L}$-geodesic, then variations
of $\gamma$ fixing the endpoints  give curves whose ${\mathcal
L}$-length is no less than that of $\gamma$. In fact, there is a
second-order version of this inequality which we shall need later.

\begin{corollary}[Positivity of the second variation]\label{cor:l-length-second-variation-positive}
\uses{prop:l-length-second-variation, lem:l-exponential-differential}
\label{posform}

Let $Z\in T_xM_{T-\tau_1}$. Suppose that the associated ${\mathcal
L}$-geodesic $\gamma_Z$ minimizes ${\mathcal L}$-length between its
endpoints, $x$ and $\gamma_Z(\bar\tau)$, and that $d_Z{\mathcal
L}\exp^{\bar\tau}_x$ is an isomorphism. Then for any family
$\gamma_u$ of curves parameterized by backward time with
$Y=(\partial \gamma/\partial u)|_{u=0}$ vanishing at both endpoints,
we have
$$\frac{d^2}{du^2}{\mathcal L}(\gamma_u)|_{u=0}\ge 0,$$
with equality if and only if $Y=0$.
\end{corollary}

\begin{proof}
\sketch
\uses{prop:l-length-second-variation,lem:l-exponential-differential}
According to Proposition~\ref{prop:l-length-second-variation} the second variation  in the
$Y$-direction is non-negative and vanishes if and only if $Y$ is an ${\mathcal
L}$-Jacobi field. But since $d_Z{\mathcal L}\exp_x^{\bar\tau}$ is a
diffeomorphism, by Lemma~\ref{lem:l-exponential-differential} there are no non-zero ${\mathcal
L}$-Jacobi fields vanishing at both endpoints of $\gamma_Z$.
\end{proof}



\subsection{The gradient of $\widetilde L^\tau$}

Recall that $\widetilde
L^{\tau}$ is the map from ${\mathcal D}^{\tau}_x$ to $\Ar$ that
assigns to each $Z$ the ${\mathcal L}$-length of
$\gamma_Z|_{[\tau_1,\tau]}$. We compute its gradient.




\begin{lemma}[Gradient of $\widetilde L^\tau$]\label{lem:l-length-tilde-gradient}
\uses{def:l-exponential-map, def:l-length}
\label{DtildeL}

Suppose that $Z\in {\mathcal D}^{\tau}_x$. Then for any $\widetilde
Y\in T_xM_{T-\tau_1}=T_Z({\mathcal D}^{\tau}_x)$ we have
$$ \langle \nabla\widetilde L^{\tau},\widetilde Y\rangle =2\sqrt{\tau}\langle
X(\tau),d_Z\left({\mathcal L}\exp_x^{\tau}\right)(\widetilde
Y)\rangle.$$
\end{lemma}


\begin{proof}
\sketch
\uses{lem:l-exponential-differential}
Since ${\mathcal D}^{\tau}_x$ is an open subset of $T_x(M_{T-\tau_1})$, it
follows that for any $\widetilde Y\in T_x(M_{T-\tau_1})$ there is a
one-parameter family $\gamma_u(\tau')=\gamma(\tau',u)$ of
 ${\mathcal L}$-geodesics, defined for $\tau_1\le \tau'\le \tau$,
starting at $x$ with $\gamma(\cdot,0)=\gamma_Z$ and with
$\frac{\partial}{\partial u}\left(\sqrt{\tau_1}X(\tau_1)\right)=\widetilde Y$.
(Again, when $\tau_1=0$, this equation is interpreted to mean
$\frac{\partial}{\partial u}\lim_{\tau'\rightarrow
0}(\sqrt{\tau'}X(\tau',u))=\widetilde Y$.) Let
$Y(\tau')=\frac{\partial}{\partial u}(\gamma(\tau',u))|_{u=0}$ be the
corresponding ${\mathcal L}$-Jacobi field along $\gamma_Z$. Since
$\gamma(\tau_1,u)=x$ for all $u$, we have $Y(\tau_1)=0$. Since
$\gamma(\cdot,u)$ is an ${\mathcal L}$-geodesic for all $u$, according to
Equation~(\ref{variation}), and in the case $\tau_1=0$, using the fact that
$\sqrt{\tau}X(\tau')$ approaches a finite limit as $\tau\rightarrow 0$, we have
$$\frac{d}{du}{\mathcal
L}(\gamma_u)|_{u=0}=2\sqrt{\tau}\langle X(\tau),Y(\tau)\rangle.$$ By
Lemma~\ref{lem:l-exponential-differential} we have $Y(\tau)=d_Z{\mathcal L}\exp_x^{\tau}(\tilde
Y)$. Thus,
$$\langle\nabla\widetilde L^{\tau},\widetilde Y\rangle=\frac{d}{du}{\mathcal
L}(\gamma_u)|_{u=0}=2\sqrt{\tau}\langle
X(\tau),Y(\tau)\rangle=2\sqrt{\tau}\langle X(\tau),d_Z({\mathcal
L}\exp^{\tau}_x)(\widetilde Y)\rangle.$$
\end{proof}


\subsection{Local diffeomorphism near the initial $\tau$}

Now let us use the nature of the ${\mathcal L}$-Jacobi equation to
study ${\mathcal L}\exp_x$ for $\tau>\tau_1$ but $\tau$
sufficiently close to $\tau_1$.

\begin{lemma}[Local diffeomorphism near the initial time]\label{lem:l-exponential-local-diffeomorphism}
\uses{lem:l-exponential-differential, def:l-jacobi-field}
\label{locstr}

For any $x$ in ${\mathcal M}$ with ${\bf
t}(x)=T-\tau_1$ and any $Z\in T_xM_{t-\tau_1}$, there is $\delta>0$
such that for any $\tau$ with $\tau_1<\tau<\tau_1+\delta$ the map
${\mathcal L}\exp_x^{\tau}$ is a local diffeomorphism from a
neighborhood of $Z$ in $T_xM_{T-\tau_1}$ to $M_{T-\tau}$.
\end{lemma}

\begin{proof}
\sketch
\uses{lem:l-exponential-differential}
Fix $x$ and $Z$ as in the statement of the lemma. To establish the result it
suffices to prove that there is $\delta>0$ such that $d_Z{\mathcal L}\mathrm{exp}_x^{\tau}$ is an isomorphism for all $\tau_1<\tau<\tau_1+\delta$. By
Lemma~\ref{lem:l-exponential-differential} it is enough to find a $\delta>0$ such that any ${\mathcal
L}$-Jacobi field $Y$ along
 $\gamma_Z$ with $\sqrt{\tau_1}\nabla_XY(\tau_1)\not= 0$ does
 not vanish on the interval $(\tau_1,\tau_1+\delta)$.
 Because the ${\mathcal L}$-Jacobi equation is linear, it suffices to consider the
 case of ${\mathcal L}$-Jacobi fields with
 $|\nabla_XY(\tau_1)|=1$. The space of such fields is identified with the
 unit sphere in $T_xM_{T-\tau_1}$.
 Let us consider first the case when $\tau_1\not= 0$. Then for any
 such tangent vector $\nabla_XY(\tau_1)\not=0$. Since $Y(\tau_1)=0$,
 it follows that $Y(\tau)\not=0$ in some interval
 $(\tau_1,\tau_1+\delta)$, where $\delta$ can depend on $Y$.  Using
 compactness of the unit sphere in the tangent space, we see that there is $\delta>0$
 independent of $Y$ so that the above holds.

 In case when $\tau_1=0$, it is convenient to shift to the
 $s=\sqrt{\tau}$ parameterization. Then the geodesic equation and
 the ${\mathcal L}$-Jacobi equation are non-singular at the origin. Also, letting
 $A=d\gamma_Z/ds$ we have $\nabla_AY=2\lim_{\tau\rightarrow
 0}\sqrt{\tau}\nabla_XY$. In these variables, the argument for $\tau_1=0$
 is the same as the one above for $\tau_1>0$.
\end{proof}

\begin{remark}[${\mathcal L}\exp$ for $\tau<\tau_1$]\label{rem:l-exponential-backward-tau}
\uses{lem:l-exponential-local-diffeomorphism, lem:l-length-tilde-gradient}
\label{backwards}

When $\tau_1>0$ it is possible to consider the ${\mathcal L}\mathrm{exp}^{\tau}_x$ defined for $0<\tau<\tau_1$. In this case, the curves
are moving backward in $\tau$ and hence are moving forward with
respect to the time parameter ${\bf t}$. Two comments are in order.
First of all, for $\tau<\tau_1$, the gradient of ${\widetilde
L}^{\tau}_x$ is $-2\sqrt{\tau}X(\tau)$. The reason for the sign
reversal is that the length is given by the integral from $\tau$ to
$\tau_1$ and hence its derivative with respect to $\tau$ is the
negative of the integrand. The second thing to remark is that
Lemma~\ref{lem:l-exponential-local-diffeomorphism} is true for $\tau<\tau_1$ with $\tau$
sufficiently close to $\tau_1$.
\end{remark}

\section{Minimizing ${\mathcal L}$-geodesics and the injectivity
domain}\label{sect:InjDom}

Now we discuss the analogue of the interior of the cut locus for the
usual exponential map of a Riemannian manifold. For
Section~\ref{sect:InjDom} we keep the assumption that $x\in
{\mathcal M}$ with ${\bf t}(x)=T-\tau_1$ for some $\tau_1\ge 0$.


\begin{definition}[Injectivity set]\label{def:l-injectivity-set}
\uses{def:l-exponential-map, def:l-length}
\label{injdefn}

 The {\sl injectivity set}
$\widetilde{\mathcal U}_x\subset {\mathcal D}_x\subset
\left(T_xM_{T-\tau_1}\times (\tau_1,\infty)\right)$ is the subset of all $(Z,\tau)\in {\mathcal D}_x$ with the following
properties:
\begin{enumerate}
\item The map ${\mathcal L}\exp_x^{\tau}$ is a local
diffeomorphism near $Z$ from $T_x(M_{T-\tau_1})$ to $M_{T-\tau}$.
\item There is a neighborhood ${\mathcal Z}$ of $Z$ in ${\mathcal D}^{\tau}_x$
 such that for every $Z'\in {\mathcal Z}$ the ${\mathcal L}$-geodesic
$\gamma_{Z'}|_{[\tau_1,\tau]}$ is the unique minimizing path
parameterized by backward time for the ${\mathcal L}$-length. That
is to say, the ${\mathcal L}$-length of
$\gamma_{Z'}|_{[\tau_1,\tau]}$ is less than the ${\mathcal
L}$-length of any other path parameterized by backward time between
the same endpoints.
\end{enumerate}
For any $\tau>\tau_1$, we set $\widetilde {\mathcal
U}_x(\tau)\subset T_xM_{T-\tau_1}$ equal to the slice of $\widetilde{\mathcal U}_x$ at
$\tau$, i.e., $\widetilde {\mathcal U}_x(\tau)$ is determined by the
equation
$$\widetilde{\mathcal
U}_x(\tau)\times\{\tau\}=\widetilde{\mathcal
U}_x\cap\left(T_xM_{T-\tau_1}\times\{\tau\}\right).$$
\end{definition}

It is clear from the definition that $\widetilde{\mathcal
U}_x\subset {\mathcal D}_x$ is an open subset and hence
$\widetilde{\mathcal U}_x$ is an open subset of
$T_xM_{T-\tau_1}\times (\tau_1,\infty)$. Of course, this implies
that $\widetilde {\mathcal U}_x(\tau)$ is an open subset of
${\mathcal D}_x^\tau$ for every $\tau>\tau_1$.



\begin{definition}[Injectivity domain]\label{def:l-injectivity-domain}
\uses{def:l-injectivity-set, def:l-exponential-map}
\label{L_x}

We set ${\mathcal U}_x\subset {\mathcal M}$ equal to $ {\mathcal L}\exp_x(\widetilde{\mathcal
U}_x)$. We call this subset of ${\mathcal M}$ the {\em injectivity
domain (of $x$)}.
 For any
$\tau>\tau_1$ we set ${\mathcal U}_x(\tau)={\mathcal U}_x\cap
M_{T-\tau}$.
\end{definition}

By definition, for every point $q\in {\mathcal U}_x$ for any
$(Z,\tau)\in \widetilde{\mathcal U}_x$ with ${\mathcal L}\mathrm{exp}_x(Z,\tau)=q$, the ${\mathcal L}$-geodesic
$\gamma_{Z}|_{[\tau_1,\tau]}$  is a minimizing
 ${\mathcal L}$-geodesic to $q$. In particular, there is a minimizing
 ${\mathcal L}$-geodesic from $x$ to $q$.

\begin{definition}[The minimal ${\mathcal L}$-length function $L_x$]\label{def:minimal-l-length-function}
\uses{def:l-injectivity-domain, def:l-length}

The function $L_x\colon {\mathcal U}_x\to \Ar$
assigns to each $q$ in ${\mathcal U}_x$ the length of any minimizing
${\mathcal L}$-geodesic from $x$ to $q$. For any $\tau>\tau_1$, we
denote by $L_x^\tau$ the restriction of $L_x$
to the $T-\tau$ time-slice of ${\mathcal U}_x$, i.e., the
restriction of $L_x$ to ${\mathcal U}_x(\tau)$.
\end{definition}

 This brings us to the
analogue of the fact that in Riemannian geometry the restriction to
the interior of the cut locus of the exponential mapping is a
diffeomorphism onto an open subset of the manifold.


\begin{proposition}[${\mathcal L}\exp_x$ is a diffeomorphism onto the injectivity domain]\label{prop:l-exponential-diffeomorphism}
\uses{def:l-injectivity-set, def:l-injectivity-domain, def:minimal-l-length-function}
\label{diffeo}

The map
$${\mathcal L}\exp_x\colon\widetilde {\mathcal
U}_x\to {\mathcal M}$$ is a diffeomorphism onto the open subset
${\mathcal U}_x$ of ${\mathcal M}$. The function $L\colon {\mathcal
U}_x\to \Ar$ that associates to each $q\in {\mathcal U}_x$ the
length of the unique minimizing ${\mathcal L}$-geodesic from $x$ to
$q$ is a smooth function and
$$L_x\circ {\mathcal L}\exp_x|_{\widetilde{\mathcal U}_x}
=\widetilde L|_{\widetilde{\mathcal U}_x}.$$
\end{proposition}



\begin{proof}
\sketch
\uses{def:l-injectivity-set}
We consider the differential of ${\mathcal L}\exp_x$ at any
 $(Z,\tau)\in\widetilde {\mathcal U}_x$. By construction
 the restriction of this differential to $T_xM_{T-\tau_1}$ is a local
 isomorphism onto ${\mathcal H}T{\mathcal M}$ at the image point. On the other
 hand, the differential of ${\mathcal L}\exp_x$ in the $\tau$
 direction is $\gamma_Z'(\tau)$, whose `vertical' component is
 $-\chi$. By the inverse function theorem this
 shows that ${\mathcal L}\exp_x$ is a local diffeomorphism at
 $(Z,\tau)$, and
 its image is an open subset of ${\mathcal M}$.
The uniqueness in Condition 2, of the definition immediately implies
that the restriction of ${\mathcal L}\exp_x$ to
$\widetilde{\mathcal U}_x$ is one-to-one, and hence that it is a
global diffeomorphism onto its image ${\mathcal U}_x$.

Since for every $(Z,\tau)\in\widetilde{\mathcal U}_x$ the ${\mathcal
L}$-geodesic  $\gamma_Z|_{[\tau_1,\tau]}$ is ${\mathcal
L}$-minimizing, we see that $L_x\circ {\mathcal L}\mathrm{exp}_x|_{\widetilde{\mathcal U}_x}=\widetilde L|_{\widetilde
{\mathcal U}_x}$ and that $L_x\colon {\mathcal U}_x\to \Ar$ is a
smooth function.
\end{proof}

According to Lemma~\ref{lem:l-length-tilde-gradient} we have:

\begin{corollary}[Gradient of $L_x^\tau$]\label{cor:l-length-gradient}
\uses{lem:l-length-tilde-gradient, prop:l-exponential-diffeomorphism}
\label{DL}

At any $q\in {\mathcal U}_x(\tau)$ we have
$$\nabla L_x^{\tau}(q)=2\sqrt{\tau}X(\tau)$$
where  $X(\tau)$ is the horizontal component of $\gamma'(\tau)$,
where $\gamma$ is the unique minimizing ${\mathcal L}$-geodesic
connecting $x$  to $q$.
\end{corollary}

At the level of generality that we are working (arbitrary
generalized Ricci flows) there is no result analogous to the fact in
Riemannian geometry that the image under the exponential mapping of
the interior of the cut locus is an open dense subset of the
manifold. There is an analogue in the special case of Ricci flows on
compact manifolds or on complete manifolds of bounded curvature.
These will be discussed in Section~\ref{newcomp2}.





\subsection{Monotonicity of the $\widetilde{\mathcal U}_x(\tau)$ with respect to
$\tau$}


Next, we have the analogue of the fact in Riemannian geometry that
the cut locus is star-shaped.


\begin{proposition}[Monotonicity of the injectivity set]\label{prop:injectivity-set-monotone}
\uses{def:l-injectivity-set}
\label{inclusion}

Let $\bar\tau'>\bar\tau$. Then $\widetilde{\mathcal
U}_x(\bar\tau')\subset \widetilde{\mathcal U}_x(\bar\tau)\subset
T_xM_{T-\tau_1}$.
\end{proposition}


\begin{proof}
\sketch
\uses{lem:l-exponential-local-diffeomorphism,lem:l-exponential-differential,prop:l-length-second-variation}
For $Z\in \widetilde{\mathcal U}_x(\bar\tau')$,
 we shall show that:
(i) the ${\mathcal L}$-geodesic $\gamma_{Z'}|_{[\tau_1,\bar\tau]}$
is the unique minimizing ${\mathcal L}$-geodesic from $x$ to
$\gamma_{Z}(\bar\tau)$, and (ii) the differential $d_Z{\mathcal
L}\exp_x^{\bar\tau}$ is an isomorphism. Given these two
conditions, it follows from the definition that $\widetilde{\mathcal
U}_x(\bar\tau')$ is contained in $\widetilde{\mathcal
U}_x(\bar\tau)$.

 We  show that the ${\mathcal L}$-geodesic
$\gamma_{Z}|_{[\tau_1,\bar\tau]}$ is  the unique minimizing
${\mathcal L}$- geodesic to its endpoint.  If there is an ${\mathcal
L}$-geodesic $\gamma_1$, distinct from
$\gamma_{Z}|_{[\tau_1,\bar\tau]}$, from $x$ to $\gamma_Z(\bar\tau)$
whose ${\mathcal L}$-length is at most that of
$\gamma_Z|_{[\tau_1,\bar\tau]}$, then we can create a broken path
$\gamma_1*\gamma_Z|_{[\bar\tau,\bar\tau']}$ parameterized by
backward time whose ${\mathcal L}$-length is at most that of
$\gamma_{Z}$. Since this latter path is not smooth, its ${\mathcal
L}$-length cannot be the minimum, which is a contradiction.

 Now suppose  that $d_Z{\mathcal L}\exp_x^{\bar\tau}$ is not an isomorphism.
The argument is similar to the one above, using a non-zero
${\mathcal L}$-Jacobi field vanishing at both endpoints rather than
another geodesic. Let $\tau'_2$ be the first $\tau$ for which
$d_Z{\mathcal L}\mathrm{exp}_x^\tau$ is not an isomorphism.
According to Lemma~\ref{lem:l-exponential-local-diffeomorphism}, $\tau_1<\tau'_2\le \bar\tau$. Since
${\mathcal L}\exp_x^{\tau'_2}$ is not a local diffeomorphism at
$(Z,\tau'_2)$, by Lemma~\ref{lem:l-exponential-differential}  there is a non-zero
${\mathcal L}$-Jacobi field $Y$ along $\gamma_Z|_{[\tau_1,\tau'_2]}$
vanishing at both ends. Since $\gamma_Z|_{[\tau_1,\tau_2']}$ is
${\mathcal L}$-minimizing, according to Proposition~\ref{prop:l-length-second-variation},
the second variation of the length of $\gamma_Z|_{\tau_1,\tau'_2]}$
in the $Y$-direction vanishes, in the sense that if $\gamma(u,\tau)$
is any one-parameter family of paths parameterized by backward time
from $x$ to $\gamma_Z(\tau'_2)$ with $(\partial \gamma/\partial
u)|_{u=0}=Y$ then
$$\frac{\partial^2{\mathcal L}(\gamma_u)}{\partial
u^2}\bigl|_{u=0}\bigr.=0.$$ Extend $Y$ to a horizontal vector field
$\widehat Y$ along $\gamma_Z$ by setting $\widehat Y(\tau)=0$ for
all $\tau\in [\tau'_2,\bar\tau]$. Of course, the extended horizontal
vector field $\widehat Y$ is not $C^2$ at $\tau'_2$ since $Y$, being
a non-zero ${\mathcal L}$-Jacobi field, does not vanish to second
order there. This is the first-order variation of the family
$\hat\gamma(u,\tau)$ that agrees with $\gamma(u,\tau)$ for all
$\tau\le \tau'_2$ and has $\hat\gamma(u,\tau)=\gamma_Z(\tau)$ for
all $\tau\in [\tau'_2,\bar\tau]$. Of course, the second-order
variation of this extended family at $u=0$ agrees with the
second-order variation of the original family at $u=0$, and hence
vanishes. But according to Proposition~\ref{prop:l-length-second-variation} this means that
$\widehat Y$ is an ${\mathcal L}$-Jacobi field, which is absurd
since it is not a $C^2$-vector field.
\end{proof}



We shall also need a closely related result.

\begin{proposition}[Local diffeomorphism of ${\mathcal L}\exp$ along a minimizing geodesic]\label{prop:l-exponential-diffeomorphism-along-geodesic}
\uses{def:l-exponential-map, def:l-injectivity-domain, def:l-geodesic-minimizing}
\label{tausmooth}

Let $\gamma$ be a minimizing ${\mathcal L}$-geodesic defined for
$[\tau_1,\bar\tau]$. Fix $0\le\tau_1<\tau_2<\bar\tau$, and set
$q_2=\gamma(\tau_2)$, and $Z_2=\sqrt{\tau_2}X_{\gamma}(\tau_2)$.
Then, the map ${\mathcal L}\exp_{q_2}$ is  diffeomorphism from
a neighborhood of $\{Z_2\}\times (\tau_2,\bar\tau]$ in
$T_{q}M_{T-\tau_2}\times (\tau_2,\infty)$ onto a neighborhood of the
image of  $\gamma|_{(\tau_2,\bar\tau]}$.
\end{proposition}

\begin{proof}
\sketch
\uses{def:l-jacobi-field,prop:l-length-second-variation}
It suffices to show that the differential of ${\mathcal L}\mathrm{exp}_{q_2}^\tau$ is an isomorphism for all $\tau\in
(\tau_2,\bar\tau]$. If this is not the case, then there is a
$\tau'\in (\tau_2,\bar\tau]$ and a non-zero ${\mathcal L}$-Jacobi
field $Y$ along $\gamma_Z|_{[\tau_2,\tau']}$ vanishing at both ends.
We extend $Y$ to a horizontal vector field $\widehat Y$ along all of
$\gamma_Z|_{[\tau_1,\tau']}$  by setting it equal to zero on
$[\tau_1,\tau_2]$. Since $Y$ is an ${\mathcal L}$-Jacobi field, the
second-order variation of ${\mathcal L}$-length in the direction of
$Y$ is zero, and consequently the second-order variation of the
length of $\gamma_Z|_{[\tau_1,\tau']}$ vanishes. Hence by
Proposition~\ref{prop:l-length-second-variation} it must be the case that $\widehat Y$ is a
${\mathcal L}$-Jacobi field. This is impossible since $\widehat Y$
is not smooth at $\tau'$.
\end{proof}



We finish this section with a computation of the $\tau$-derivative
of $L_x$.

\begin{lemma}[$\tau$-derivative of $L_x$]\label{lem:l-length-tau-derivative}
\uses{def:minimal-l-length-function, def:l-injectivity-domain}
\label{Ltau}

 Suppose that
$q\in  {\mathcal U}_x$ with ${\bf t}(q)=T-\bar\tau$ for some $\bar\tau>\tau_1$.
Let $\gamma\colon [\tau_1,\bar\tau]\to {\mathcal M}$ be the unique minimizing
${\mathcal L}$-geodesic from $x$ to $q$. Then we have
\begin{equation}\label{Ltaueqn}
\frac{\partial L_x}{\partial \tau}(q)=2\sqrt{\bar\tau}R(
q)-\sqrt{\bar\tau}\left(R( q)+|X(\bar\tau)|^2\right).
\end{equation}
\end{lemma}


\begin{proof}
\sketch
\uses{cor:l-length-gradient}
By definition and the Fundamental Theorem of Calculus, we have
$$\frac{d}{d\tau}L_x(\gamma(\tau))= \sqrt{\tau}\left(R(\gamma(\tau))+|X(\tau)|^2\right).$$
On the other hand since $\gamma'(\tau)=-\partial/\partial t+X(\tau)$
the chain rule implies
$$\frac{d}{d\tau}L_x(\gamma(\tau))=\langle \nabla
L_x,X(\tau)\rangle+\frac{\partial L_x}{\partial
\tau}(\gamma(\tau)),$$ so that
$$\frac{\partial L_x}{\partial \tau}(\gamma(\tau))=\sqrt{\tau}\left(R(\gamma(\tau))+|X(\tau)|^2\right)-\langle\nabla
L_x,X(\tau)\rangle.$$ Now using Corollary~\ref{cor:l-length-gradient}, and rearranging
the terms gives the result.
\end{proof}



\section{Second-order differential inequalities for $\widetilde L^{\bar\tau}$
and $L_x^{\bar\tau}$}\label{secondorder}

Throughout Section~\ref{secondorder} we fix $x\in {\mathcal M}$ with $x\in
M_{T-\tau_1}$.


\subsection{The second variation formula for $\widetilde L^{\bar\tau}$}

Our goal here is to compute the second variation of $\widetilde
L^{\bar\tau}$ in the direction of a horizontal vector field
$Y(\tau)$ along an ${\mathcal L}$-geodesic $\gamma$. Here is the
main result of this subsection.

\begin{proposition}[Second variation formula for $\widetilde L^{\bar\tau}$]\label{prop:l-length-second-variation-general}
\uses{def:l-geodesic, def:l-length}
\label{2ndvariation}

 Fix $0\le \tau_1<\bar \tau$.
Let $\gamma$ be an ${\mathcal L}$-geodesic defined on
$[\tau_1,\bar\tau]$ and let $\gamma_u=\widetilde\gamma(\tau,u)$ be a
smooth family of curves parameterized by backward time with
$\gamma_0=\gamma$. Let $\widetilde Y(\tau,u)$ be $\partial
\widetilde\gamma/\partial u$ and let $\widetilde X$ be the
horizontal component of $\partial \widetilde \gamma/\partial \tau$.
These are horizontal vector fields along the image of $\widetilde
\gamma$. We set $Y$ and $X$ equal to the restrictions of $\widetilde
Y$ and $\widetilde X$, respectively, to $\gamma$. We assume that
$Y(\tau_1)=0$. Then
\begin{eqnarray*}
\lefteqn{\frac{d^2}{du^2}\left({\mathcal L}(\gamma_u)\right)|_{u=0}
 =  2\sqrt{\bar\tau}\langle \nabla_{Y(\tau)}\widetilde
Y(\bar\tau,u)|_{u=0},X(\bar\tau)\rangle } \\
& & + \int_{\tau_1}^{\bar\tau}\sqrt{\tau}(\Hess(R)(Y,Y)
        +2\langle{\mathcal R}(Y,X)Y,X\rangle
        - 4(\nabla_{Y}\Ric)(X,Y) \\
        & &  + 2(\nabla_{X}\Ric)(Y,Y)+ 2\abs{\nabla_{X} Y}^{2})d\tau.
\end{eqnarray*}
\end{proposition}


As we shall see, this is simply a rewriting of the equation in
Lemma~\ref{lem:l-length-second-variation-formula} in the special case when $u_1=u_2$.



We begin the proof of this result with the following computation.

\begin{claim}[Rate of change of $\langle \nabla_X Y,Y\rangle$ along $\tau$]\label{claim:covariant-derivative-inner-product-formula}
\uses{def:ricci-curvature}

  Let
$\gamma(\tau)$ be a curve parameterized by backward time. Let $Y$ be
a horizontal vector field along $\gamma$ and let $X$ be the
horizontal component of $\partial\widetilde\gamma/\partial \tau$.
Then
\begin{align*}\frac{\partial}{\partial\tau}\langle \nabla_{ X}Y, Y\rangle
 &= \langle \nabla_{ X} Y, \nabla_{ X} Y\rangle  + \langle
\nabla_{ X}\nabla_{ X} Y, Y\rangle
\\&\hspace{0.5cm} +2\Ric(\nabla_XY,Y)) +(\nabla_X\Ric)(Y,Y))\end{align*}
\end{claim}

\begin{proof}
  \uses{lem:l-length-second-variation-formula, lem:l-jacobi-equation, claim:covariant-derivative-inner-product-formula}
\sketch
We can break $\frac{\partial}{\partial\tau}\langle \nabla_{X}Y,Y\rangle $ into
two parts: the first  assumes that the metric is constant and the second deals
with the variation with $\tau$ of the metric. The first contribution is the
usual formula
$$\frac{\partial}{\partial\tau}\langle \nabla_{X}Y,Y\rangle _{G(T-\tau_0)}=
\langle \nabla_{X}Y,\nabla_{X}Y\rangle _{G(T-\tau_0)}+ \langle
\nabla_{X}\nabla_{X}Y,Y\rangle _{G(T-\tau_0)}.$$ This gives us the first two
terms of the right-hand side of the equation in the claim.

We show that the last two terms in that equation come from differentiating the
metric with respect to $\tau$. To do this recall that in local coordinates,
writing the metric $G(T-\tau)$ as $g_{ij}$, we have
$$\langle \nabla_{X}Y,Y \rangle=g_{ij}\bigl(X^k\partial_k
Y^i+\Gamma_{kl}^iX^kY^l\bigr)Y^j.$$ There are two contributions coming from
differentiating the metric with respect to $\tau$. The first is when we
differentiate $g_{ij}$. This leads to
$$2\Ric_{ij}\bigl(X^k\partial_k
Y^i+\Gamma_{kl}^iX^kY^l\bigr)Y^j=2\Ric(\nabla_XY,Y\rangle.$$
The other contribution is from differentiating the Christoffel
symbols. This yields
$$g_{ij}\frac{\partial\Gamma^{i}_{kl}}{\partial \tau}X^{k}Y^{l}Y^{j}.$$
Differentiating the formula
$\Gamma_{kl}^i=\frac{1}{2}g^{si}(\partial_kg_{sl}+\partial_lg_{sk}-\partial_s
g_{kl})$ leads to \begin{eqnarray*} g_{ij}\frac{\partial
\Gamma_{kl}^i}{\partial \tau} & = & -2\Ric_{ij}\Gamma_{kl}^i+g_{ij}g^{si}(\partial_k\Ric_{sl}+\partial_l\Ric_{sk}
-\partial_s\Ric_{kl})\\
& = & -2\Ric_{ij}\Gamma_{kl}^i+
\partial_k\Ric_{jl}+\partial_l\Ric_{jk} -\partial_j\Ric_{kl}.
\end{eqnarray*}
Thus, we have
\begin{eqnarray*}
g_{ij}\frac{\partial\Gamma^{i}_{kl}}{\partial \tau}X^{k}Y^{l}Y^{j}
 & = & \bigl(-2\Ric_{ij}\Gamma^i_{kl}+\partial_k\Ric_{jl})\bigr)X^kY^lY^j \\
& = & (\nabla_X\Ric)(Y,Y)
\end{eqnarray*}
This completes the proof of the claim.
\end{proof}








Now we return to the proof of the second variational formula in
Proposition~\ref{prop:l-length-second-variation-general}.


\begin{proof}
\sketch
\uses{lem:l-length-second-variation-formula,lem:l-jacobi-equation,claim:covariant-derivative-inner-product-formula}
According to Lemma~\ref{lem:l-length-second-variation-formula} we have
$$\frac{d^2}{du^2}{\mathcal L}_{u=0}=2\sqrt{\bar
\tau}Y(\bar\tau)(\langle \widetilde Y(\bar\tau,u),\widetilde
X(\bar\tau,u)\rangle)|_{u=0}-\int_{\tau_1}^{\tau_2}2\sqrt{\tau}\langle
\Jac(Y),Y\rangle d\tau.$$ We plug in
Equation~\ref{jaceqn} for $\Jac(Y)$ and this makes the
integrand
\begin{eqnarray*}\sqrt{\tau}\langle\nabla_Y(\nabla
R),Y\rangle+2\sqrt{\tau}\langle{\mathcal R}(Y,X)Y,X\rangle
-\bigl(2\sqrt{\tau}\langle\nabla_X\nabla_XY,Y\rangle+\frac{1}{\sqrt{\tau}}\langle\nabla_XY,Y\rangle\bigr)\\
-4\sqrt{\tau}(\nabla_Y\Ric)(X,Y)-4\sqrt{\tau}\Ric(\nabla_XY,Y)
\end{eqnarray*}
The first term is $\sqrt{\tau}\Hess(R)(Y,Y)$. Let us deal with
the third and fourth terms, which are grouped together within
parentheses. According to the previous claim, we have
\begin{eqnarray*}
\frac{\partial}{\partial \tau}\bigl(2\sqrt{\tau}\langle
\nabla_XY,Y\rangle\bigr) & = &
\frac{1}{\sqrt{\tau}}\langle\nabla_XY,Y\rangle+2\sqrt{\tau}\langle\nabla_X\nabla_XY,Y\rangle
 +2\sqrt{\tau}\langle\nabla_XY,\nabla_XY\rangle \\ & & +4\sqrt{\tau}\Ric(\nabla_XY,Y)
 +2\sqrt{\tau} (\nabla_X\Ric)(Y,Y).
\end{eqnarray*}
This allows us to replace the two terms under consideration by
$$-\frac{\partial}{\partial t}\bigl(2\sqrt{\tau}\langle
\nabla_XY,Y\rangle\bigr)+2\sqrt{\tau}\langle\nabla_XY,\nabla_XY\rangle+4\sqrt{\tau}\Ric(\nabla_XY,Y) +2\sqrt{\tau} (\nabla_X\Ric)(Y,Y).$$
Integrating the total derivative out of the integrand  and canceling
terms leaves the integrand as
\begin{eqnarray*}\sqrt{\tau}\Hess(R)(Y,Y)+2\sqrt{\tau}\langle{\mathcal R}(Y,X)Y,X\rangle
+2\sqrt{\tau}|\nabla_XY|^2 \\
-4\sqrt{\tau}(\nabla_Y\Ric)(X,Y)+2\sqrt{\tau} (\nabla_X\Ric)(Y,Y),
\end{eqnarray*}
and makes the boundary term (the one in front of the integral) equal
to
$$2\sqrt{\bar\tau}\bigl(Y(\bar\tau)\langle
\widetilde Y(\bar\tau,u),\widetilde
X(\bar\tau,u)\rangle|_{u=0}-\langle
\nabla_XY(\bar\tau),Y(\bar\tau)\rangle\bigr)=2\sqrt{\bar\tau}\langle
X(\bar\tau),\nabla_Y\widetilde Y(\bar\tau,u)|_{u=0}\rangle.$$ This
completes the proof of the proposition.
\end{proof}








\subsection{Inequalities for the Hessian of $L_x^{\bar\tau}$}

Now we shall specialize the type of vector fields along $\gamma$.
This will allow us to give an inequality for the Hessian of
${\mathcal L}$ involving the integral of the vector field along
$\gamma$. These lead to inequalities for the Hessian of
$L_x^{\bar\tau}$. The main result
of this section is Proposition~\ref{prop:l-length-hessian-inequality} below. In the end we
are interested in the case when the $\tau_1=0$. In this case the
formulas simplify. The reason for working here in the full
generality of all $\tau_1$ is in order to establish differential
inequalities at points not in the injectivity domain. As in the case
of the theory of geodesics, the method is to establish weak
inequalities at these points by working with smooth barrier
functions. In the geodesic case the barriers are constructed by
moving the initial point out the geodesic a small amount. Here the
analogue is to move the initial point of an ${\mathcal L}$-geodesic
from $\tau_1=0$ to a small positive $\tau_1$. Thus, the case of
general $\tau_1$ is needed so that we can establish the differential
inequalities for the barrier functions that yield the weak
inequalities at non-smooth points.

\begin{definition}[Adapted vector field]\label{def:adapted-vector-field}
\uses{def:l-injectivity-domain, def:ricci-curvature}

Let $q\in {\mathcal U}_x(\bar\tau)$ and let $\gamma\colon
[\tau_1,\bar\tau]\to {\mathcal M}$ be the unique minimizing
${\mathcal L}$-geodesic from $x$ to $q$. We say that a horizontal
vector field $Y(\tau)$ along $\gamma$ is {\em adapted} if it solves
the following ODE on $[\tau_1,\bar \tau]$:
\begin{equation}\label{firstYeqn}
\nabla_{X}Y(\tau) = -\Ric(Y(\tau),\cdot)^* +
\frac{1}{2\sqrt{\tau}(\sqrt{\tau}-\sqrt{\tau_1})}Y(\tau).
\end{equation}
 \end{definition}

Direct computation shows the following:

 \begin{lemma}[Norm of an adapted vector field]\label{lem:adapted-vector-field-norm}
\uses{def:adapted-vector-field}
\label{adaptform}
 Suppose that $Y(\tau)$ is an adapted vector field along $\gamma$.
 Then
\begin{eqnarray}\label{secondYeqn}
\frac{d}{d\tau}\langle Y(\tau),Y(\tau)\rangle & = &
2\Ric(Y(\tau),Y(\tau)) + 2\langle \nabla_{X}Y(\tau),Y(\tau)\rangle \\
& = &  \frac{1}{\sqrt{\tau}(\sqrt{\tau}-\sqrt{\tau_1})}\langle
Y(\tau),Y(\tau)\rangle\nonumber.
\end{eqnarray} It follows that
$$|Y(\tau)|^2 = C\frac{(\sqrt{\tau}-\sqrt{\tau_1})^2}{(\sqrt{\bar\tau}-\sqrt{\tau_1})^2},$$
where $C=|Y(\bar\tau)|^2$.
\end{lemma}





Now we come to the main result of this subsection, which is an
extremely important inequality for the Hessian of $L_x^{\bar\tau}$.

\begin{proposition}[Hessian inequality for $L_x^{\bar\tau}$]\label{prop:l-length-hessian-inequality}
\uses{def:adapted-vector-field, def:minimal-l-length-function}
\label{Hessineq}

Suppose that $q\in {\mathcal U}_x(\bar\tau)$, that $Z\in
\widetilde{\mathcal U}_x(\bar\tau)$ is the pre-image of $q$,  and
that $\gamma_Z$ is the  ${\mathcal L}$-geodesic to $q$ determined by
$Z$. Suppose  that $Y(\tau)$ is an adapted vector field along
$\gamma_Z$. Then
\begin{equation}\label{Hessin}
\Hess(L_x^{\bar\tau})(Y(\bar\tau),Y(\bar\tau)) \leq
\left(\frac{|Y(\bar\tau)|^2}{\sqrt{\bar\tau}-\sqrt{\tau_1}}\right)-
2\sqrt{\bar\tau}\Ric(Y(\bar\tau),Y(\bar\tau)) -
\int_{\tau_1}^{\bar\tau}\sqrt{\tau}H(X,Y)d\tau,
\end{equation} where
\begin{eqnarray} \nonumber
H(X,Y) & =&-\Hess(R)(Y,Y) - 2\langle  {\mathcal R}(Y,X)Y,X\rangle
\\& &-4(\nabla_{X}\Ric)(Y,Y) +4(\nabla_{Y}\Ric)(Y,X)\label{Heqn}
\\& & - 2\frac{\partial\Ric}{\partial\tau}(Y,Y) + 2\abs{\Ric(Y,\cdot)}^{2} -
\frac{1}{\tau}\Ric(Y,Y),\nonumber.
\end{eqnarray}
 We
have equality in Equation~(\ref{Hessin}) if and only if the adapted vector
field $Y$ is also a ${\mathcal L}$-Jacobi field.
\end{proposition}

\begin{remark}[$H(X,Y)$ is quadratic in $Y$]\label{rem:h-quadratic-in-y}
\uses{prop:l-length-hessian-inequality}

In spite of the notation, $H(X,Y)$ is a purely quadratic function of
the vector field $Y$ along $\gamma_Z$.
\end{remark}

We begin the proof of this proposition with three elementary lemmas. The first
is an immediate consequence of the definition of $\widetilde {\mathcal
U}_x(\bar\tau)$.


\begin{lemma}[Existence of a family of ${\mathcal L}$-geodesics with prescribed variation]\label{lem:l-geodesic-family-existence}
\uses{def:l-injectivity-domain}

Suppose that $q\in {\mathcal U}_x(\bar\tau)$  and that $\gamma\colon
[\tau_1,\bar\tau]\to{\mathcal M}$ is the minimizing ${\mathcal
L}$-geodesic from $x$ to $q$. Then for every tangent vector $
Y(\bar\tau)\in T_qM_{T-\bar\tau}$ there is a one-parameter family of
${\mathcal L}$-geodesics   $\tilde\gamma(\tau,u)$ defined on
$[\tau_1,\bar\tau]$ with $\tilde\gamma(0,u)=x$ for all $u$, with
$\tilde\gamma(\tau,0)=\gamma(\tau)$ and $\partial\tilde
\gamma(\bar\tau,0)/\partial u=Y(\bar \tau)$. Also, for every $Z\in
T_xM_{T-\tau_1}$ there is a family of ${\mathcal L}$-geodesics
$\tilde\gamma(\tau,u)$ such that $\gamma(0,u)=x$ for all $u$,
$\tilde\gamma(\tau,0)=\gamma(\tau)$ and such that, setting
$Y(\tau)=\frac{\partial}{\partial u} \tilde\gamma_u(\tau)|_{u=0}$,
we have
$$\nabla_{\sqrt{\tau_1}X(\tau_1)}Y(\tau_1)=Z.$$
\end{lemma}

\begin{lemma}[Hessian of $L_x^{\bar\tau}$ via ${\mathcal L}$-Jacobi fields]\label{lem:l-length-hessian-jacobi-formula}
\uses{def:l-jacobi-field, def:minimal-l-length-function}
\label{Jacobi}

Let  $\gamma$ be a minimizing ${\mathcal L}$-geodesic from $x$, and
let $Y(\tau)$ be an ${\mathcal L}$-Jacobi field along $\gamma$. Then
$$2\sqrt{\bar\tau}\langle\nabla_XY(\bar\tau),Y(\bar\tau)\rangle = \Hess(L_x^{\bar\tau})(Y(\bar\tau),
Y(\bar\tau)).$$
\end{lemma}

\begin{proof}
\sketch
\uses{cor:l-length-gradient}
Let $\gamma(\tau,u)$ be a one-parameter family of ${\mathcal
L}$-geodesics emanating from $x$ with $\gamma(u,0)$ being the
${\mathcal L}$-geodesic in the statement of the lemma  and with
$\frac{\partial}{\partial u}\gamma(\tau,0)=Y(\tau)$. We have the
extensions of $X(\tau)$ and $Y(\tau)$ to vector fields $\widetilde
X(\tau,u)$ and $\widetilde Y(\tau,u)$ defined at $\gamma(\tau,u)$
for all $\tau$ and $u$. Of course,
\begin{eqnarray*}
\lefteqn{2\sqrt{\bar\tau}\langle \nabla_Y\widetilde
X(\bar\tau,u)|_{u=0},Y(\bar\tau)\rangle  } \\
& = & Y(\langle 2\sqrt{\bar\tau}\widetilde X(\bar\tau,u),\widetilde
Y(\bar\tau,u)\rangle)|_{u=0}-\langle
2\sqrt{\bar\tau}X(\bar\tau),\nabla_Y\widetilde
Y(\bar\tau,u)|_{u=0}\rangle.\end{eqnarray*} Then by
Corollary~\ref{cor:l-length-gradient} we have
\begin{align*}
2\sqrt{\bar\tau}\langle \nabla_Y\widetilde
X(\bar\tau,u)|_{u=0},Y(\bar\tau)\rangle & = Y(\bar\tau)(\langle
\nabla L_x^{\bar\tau},\widetilde
Y(\bar\tau,u)\rangle)|_{u=0}-\langle \nabla
L_x^{\bar\tau},\nabla_{Y(\bar\tau)}\widetilde
Y(\bar\tau,u)|_{u=0}\rangle
\\
& = Y(\bar\tau)(\widetilde Y(\bar\tau,u)
L_x^{\bar\tau})|_{u=0}-\nabla_{Y(\bar\tau)}\widetilde
Y(\bar\tau,u)|_{u=0}(L_x^{\bar\tau}) \\ & =  \Hess(L_x^{\bar\tau})(Y(\bar\tau),Y(\bar\tau)).
\end{align*}
\end{proof}


Now suppose that we have a horizontal vector field that is both
adapted and ${\mathcal L}$-Jacobi. We get:

\begin{lemma}[Adapted ${\mathcal L}$-Jacobi field norm relation]\label{lem:adapted-jacobi-hessian-relation}
\uses{def:adapted-vector-field, def:l-jacobi-field, def:minimal-l-length-function}
\label{adapjab}

Suppose that $q\in {\mathcal U}_x(\bar\tau)$, that $Z\in \widetilde
{\mathcal U}_x(\bar\tau)$ is the pre-image of $q$, and that
$\gamma_Z$ is the  ${\mathcal L}$-geodesic to $q$ determined by $Z$.
Suppose further that $Y(\tau)$ is a horizontal vector field along
$\gamma$ that is both adapted and an ${\mathcal L}$-Jacobi field.
Then,  we have
$$\frac{1}{2\sqrt{\bar\tau}(\sqrt{\bar\tau}-\sqrt{\tau_1})}|Y(\bar\tau)|^2
=\frac{1}{2\sqrt{\bar\tau}}\Hess(L_x^{\bar\tau})(Y(\bar\tau),Y(\bar\tau))+\Ric(Y(\bar\tau),Y(\bar\tau)).$$
\end{lemma}

\begin{proof}
  \uses{lem:l-length-hessian-jacobi-formula, lem:l-geodesic-family-existence, prop:l-length-second-variation-general, cor:l-length-gradient, lem:adapted-vector-field-norm, prop:l-length-second-variation, lem:l-length-second-variation-formula}
\sketch

From the definition of an adapted vector field  $Y(\tau)$  we have
\begin{equation*}
\Ric(Y(\tau),Y(\tau)) + \langle
\nabla_{X}Y(\tau),Y(\tau)\rangle =
\frac{1}{2\sqrt{\tau}(\sqrt{\tau}-\sqrt{\tau_1})}\langle
Y(\tau),Y(\tau)\rangle.
\end{equation*}
Since $Y(\tau)$ is an ${\mathcal L}$-Jacobi field, according to
Lemma~\ref{lem:l-length-hessian-jacobi-formula} we have
$$\langle\nabla_XY(\bar\tau),Y(\bar\tau)\rangle =
\frac{1}{2\sqrt{\bar\tau}}\Hess(L_x^{\bar\tau})(Y(\bar\tau),Y(\bar\tau)).$$ Putting these
together gives the result.
\end{proof}

Now we are ready to begin the proof of Proposition~\ref{prop:l-length-hessian-inequality}.


\begin{proof}
  \uses{lem:l-geodesic-family-existence, prop:l-length-second-variation-general, cor:l-length-gradient, lem:adapted-vector-field-norm, prop:l-length-second-variation, lem:l-length-second-variation-formula}
\sketch

Let $\tilde\gamma(\tau,u)$ be a family  of curves with
$\gamma(\tau,0)=\gamma_Z$ and with $\frac{\partial}{\partial
u}\gamma(\tau,u)=\widetilde Y(\tau,u)$. We denote by $Y$ the
horizontal vector field which is the restriction of $\widetilde Y$
to $\gamma_0=\gamma_Z$. We denote by
$q(u)=\tilde\gamma(\bar\tau,u)$. By restricting to a smaller
neighborhood of $0$ in the $u$-direction, we can assume that
$q(u)\in {\mathcal U}_x(\bar\tau)$ for all $u$. Then ${\mathcal
L}(\tilde\gamma_u)\ge L_x^{\bar\tau}(q(u))$. Of course,
$L_x^{\bar\tau}(q(0))={\mathcal L}(\gamma_Z)$. This implies that
$$\frac{d}{d u}L_x^{\bar\tau}(q(u))\bigl|\bigr._{u=0}=
\frac{d}{du}{\mathcal L}(\gamma_u)\bigl|\bigr._{u=0},$$
 and
$$
Y(\bar\tau)(\widetilde
Y(\bar\tau,u)(L^{\bar\tau}_x))|_{u=0}=\frac{d^2}{du^2}L_x^{\bar\tau}(q(u))\bigl|\bigr._{u=0}\le
\frac{d^2}{du^2}{\mathcal L}(\gamma_u)\bigl|\bigr._{u=0}.$$
 Recall that
$\nabla L_x^{\bar\tau}(q) = 2\sqrt{\bar\tau}X(\bar\tau)$, so that
$$\nabla_{Y(\bar\tau)}\widetilde Y(\bar\tau,u)|_{u=0}(L_x^{\bar\tau})=
\langle \nabla_{Y(\bar\tau)}\widetilde Y(\bar\tau,u)|_{u=0},\nabla
L^{\bar\tau}\rangle = 2\sqrt{\bar\tau}\langle
\nabla_{Y(\bar\tau)}\widetilde
Y(\bar\tau,u)|_{u=0},X(\bar\tau)\rangle .$$ Thus, by
Proposition~\ref{prop:l-length-second-variation-general}, and using the fact that
$Y(\tau_1)=0$, we have
\begin{eqnarray*}
\lefteqn{\Hess(L^{\bar\tau})(Y(\bar\tau),Y(\bar\tau))  =
Y(\bar\tau)\left(\widetilde
Y(\bar\tau,u)(L_x^{\bar\tau})\right)|_{u=0}-
\nabla_{Y(\bar\tau)}\widetilde Y(\bar\tau,u)|_{u=0}(L_x^{\bar\tau})} \\
& \le & \frac{d^2}{du^2}{\mathcal
L}(\gamma_u)-2\sqrt{\bar\tau}\langle
\nabla_{Y(\bar\tau)}\widetilde Y(\bar\tau,u)|_{u=0},X(\bar\tau)\rangle \\
 & = & \int_{\tau_1}^{\bar\tau}
\sqrt{\tau}\bigl(\Hess(R)(Y,Y) +
2\langle {\mathcal R}(Y,X)Y,X\rangle - 4(\nabla_{Y}\Ric)(X,Y)\bigr. \\
& & \ \ \ \ \bigl.+ 2(\nabla_{X}\Ric)(Y,Y) +
2\abs{\nabla_XY}^2\bigr)d\tau.
\end{eqnarray*}
 Plugging in Equation~(\ref{firstYeqn}), and using the fact that
 $|Y(\tau)|^2=
|Y(\bar\tau)|^2\frac{(\sqrt{\tau}-\sqrt{\tau_1})^2}{(\sqrt{\bar\tau}-\sqrt{\tau_1})^2}$,
gives
\begin{eqnarray*}
\lefteqn{\Hess(L^{\bar\tau})(Y(\bar\tau),Y(\bar\tau)) } \\ &
\le & \int_{\tau_1}^{\bar\tau} \sqrt{\tau}\bigl(\Hess(R)(Y,Y) +
2\langle {\mathcal R}(Y,X)Y,X\rangle  - 4(\nabla_{Y}\Ric)(X,Y)\bigr. \\
& &  \ \ \ \ \bigl. +2(\nabla_X\Ric)(Y,Y) +
 2\abs{\Ric(Y,\cdot)}^{2}\bigr)d\tau \\ & &
 +\int_{\tau_1}^{\bar\tau}\left[
\frac{|Y(\bar\tau)|^2}{2\sqrt{\tau}(\sqrt{\bar\tau}-\sqrt{\tau_1})^2}
- \frac{2}{(\sqrt{\tau}-\sqrt{\tau_1})}\Ric(Y,Y)\right] d\tau
\end{eqnarray*}
 Using the definition of $H(X,Y)$ given in the statement, Equation~(\ref{Heqn}),
 allows us to write
\begin{eqnarray*}
\lefteqn{\Hess(L^{\bar\tau})(Y(\bar\tau),Y(\bar\tau))} \\
 & \le &
-\int_{\tau_1}^{\bar\tau}\sqrt{\tau}H(X,Y)d\tau \\ & &
+\int_{\tau_1}^{\bar\tau}\Bigl[\sqrt{\tau} \bigl(-2(\nabla_{X}\Ric)(Y,Y)
-2\frac{\partial\Ric}{\partial\tau}(Y,Y) +4|\Ric(Y,\cdot)|^2 \bigr) \Bigr.\\
& & \Bigl.+
\frac{|Y(\bar\tau)|^2}{2\sqrt{\tau}(\sqrt{\bar\tau}-\sqrt{\tau_1})^2}
-\left(\frac{2}{(\sqrt{\tau}-
\sqrt{\tau_1})}+\frac{1}{\sqrt{\tau}}\right) \Ric(Y,Y)\Bigr]d\tau,
\end{eqnarray*}

 To simplify
further, we compute, using Equation~(\ref{firstYeqn})
\begin{eqnarray*}
    \frac{d}{d\tau}\bigl(\Ric(Y(\tau),Y(\tau))\bigr) &= & \frac{\partial\Ric}{\partial{\tau}}(Y,Y) +
        2\Ric(\nabla_{X}Y,Y) + (\nabla_{X}\Ric)(Y,Y)\\
        & = & \frac{\partial\Ric}{\partial{\tau}}(Y,Y)+ (\nabla_{X}\Ric)(Y,Y) \\ & & +
\frac{1}{\sqrt{\tau}(\sqrt{\tau}-\sqrt{\tau_1})}\Ric(Y,Y) -
2|\Ric(Y,\cdot)|^{2}.
\end{eqnarray*}
 Consequently, we have

\begin{eqnarray*}
\frac{d\left(2\sqrt{\tau}\Ric(Y(\tau),Y(\tau))\right)}{d\tau} &
= & 2\sqrt{\tau}\left(\frac{\partial\Ric}{\partial{\tau}}(Y,Y)
+ (\nabla_{X}\Ric)(Y,Y)- 2|\Ric(Y,\cdot)|^{2}\right)\\
 & & + \left(\frac{2}{(\sqrt{\tau}-\sqrt{\tau_1})} +\frac{1}{\sqrt{\tau}}\right)\Ric(Y,Y)
\end{eqnarray*}
Using this, and the fact that $Y(\tau_1)=0$, gives
\begin{eqnarray}\label{Hessineq1}
\lefteqn{\Hess(L_x^{\bar\tau})(Y(\bar\tau),Y(\bar\tau)) \le}  \\
 &  &
-\int_{\tau_1}^{\bar\tau}\left(\sqrt{\tau}H(X,Y)
-\frac{d}{d\tau}\left(2\sqrt{\tau}\Ric(Y,Y)\right)
-\frac{|Y(\bar\tau)|^2}{2\sqrt{\tau}(\sqrt{\bar\tau}-\sqrt{\tau_1})^2}\right)d\tau \nonumber\\
 & = &
\frac{|Y(\bar\tau)|^2}{\sqrt{\bar\tau}-\sqrt{\tau_1}}
-2\sqrt{\bar\tau}\Ric(Y(\bar\tau),Y(\bar\tau))-\int_{\tau_1}^{\bar\tau}\sqrt{\tau}H(X,Y)d\tau.
\nonumber
\end{eqnarray}

This proves Inequality~(\ref{Hessin}). Now we examine when equality
holds in this expression. Given an adapted vector field $Y(\tau)$
along $\gamma$, let $\mu(v)$ be a geodesic through
$\gamma(\bar\tau,0)$ with tangent vector $Y(\bar\tau)$. Then there
is a one-parameter family $\mu(\tau,v)$ of minimizing ${\mathcal
L}$-geodesics with the property that $\mu(\bar\tau,v)=\mu(v)$. Let
$\widetilde Y'(\tau,v)$ be $\partial\mu(\tau,v)/\partial v$. It is
an ${\mathcal L}$-Jacobi field with $\widetilde
Y'(\bar\tau,0)=Y(\bar\tau)$. Since $L_x\circ {\mathcal L}\mathrm{exp}_x=\widetilde L$, we see that
$$\frac{d^2}{d v^2}{\mathcal
L}(\mu_v)|_{v=0}=\frac{d^2}{du^2}L^{\bar \tau}_x(\mu(u))|_{u=0}.$$
Hence, the assumption that we have equality in~(\ref{Hessin})
implies that
$$\frac{d^2}{d v^2}{\mathcal
L}(\mu_v)|_{v=0}=\frac{d^2}{d u^2}{\mathcal L}(\tilde\gamma_u)|_{u=0}.$$

 Now we extend this
one-parameter family to a two-parameter family $\mu(\tau,u,v)$ so
that $\partial\mu(\tau,0,0)/\partial v=\widetilde Y'$ and
$\partial\mu(\tau,0,0)/\partial u=Y(\tau)$. Let $w$ be the variable
$u-v$, and let $\widetilde W$ be the tangent vector in this
coordinate direction, so that $\widetilde W=\widetilde Y-\widetilde
Y'$. We denote by $W$ the restriction of $\widetilde W$ to
$\gamma_{0,0}=\gamma_Z$. By Remark~\ref{rem:second-variation-determined-by-jacobi} the second
partial derivative of the length of this family in the $u$-direction
at $u=v=0$ agrees with the second derivative of the length of the
original family $\widetilde \gamma$ in the $u$-direction.


If Inequality~(\ref{Hessin}) is an equality, then
$$\frac{\partial^2}{\partial v^2}{\mathcal L}(\mu)|_{u=v=0}=\frac{\partial^2}{\partial u^2}{\mathcal
L}(\mu)|_{u=v=0}.$$ We write $\partial/\partial u=\partial/\partial
v+\partial/\partial w$. Expanding out the right-hand side and
canceling the common terms gives
$$0=\left(\frac{\partial}{\partial v}\frac{\partial}{\partial w}+
\frac{\partial}{\partial w}\frac{\partial}{\partial
v}+\frac{\partial^2}{\partial w^2}\right){\mathcal
L}(\mu)|_{u=v=0}.$$ The previous claim tells us that the first two
terms on the right-hand side of this equation vanish, and hence we
conclude
$$\frac{\partial^2}{\partial w^2}{\mathcal L}(\mu)|_{u=v=0}=0$$
Since $W$ vanishes at both endpoints this implies, according to
Proposition~\ref{prop:l-length-second-variation}, that $\widetilde W(\tau,0,0)=0$ for all
$\tau$, or in other words $Y(\tau)=\widetilde Y'(\tau,0,0)$ for all
$\tau$. Of course by construction $\widetilde Y'(\tau,0,0)$ is an
${\mathcal L}$-Jacobi field. This shows that equality holds only if
the adapted vector field $Y(\tau)$ is also an ${\mathcal L}$-Jacobi
field.

Conversely, if the adapted vector field $Y(\tau)$ is also an
${\mathcal L}$-Jacobi field, then inequality between the second
variations at the beginning of the proof is an equality. In the rest
of the argument we dealt only with equalities. Hence, in this case
Inequality~(\ref{Hessin}) is an equality.

This shows that we have equality in~(\ref{Hessin}) if and only if
the adapted vector field $Y(\tau)$ is also an ${\mathcal L}$-Jacobi
field.
\end{proof}


\begin{claim}[Mixed second partials vanish]\label{claim:mixed-partials-vanish}
\uses{lem:l-length-second-variation-formula}

$$\frac{\partial}{\partial v}\frac{\partial}{\partial w}{\mathcal L}(\mu)|_{u=v=0}=
\frac{\partial}{\partial w}\frac{\partial}{\partial v}{\mathcal
L}(\mu)|_{v=w=0}=0.$$
\end{claim}
\begin{proof}
\sketch Of course, the second partial derivatives are equal. According to
Lemma~\ref{lem:l-length-second-variation-formula} we have
$$\frac{\partial}{\partial v}\frac{\partial}{\partial w}{\mathcal L}(\mu)|_{v=w=0}
=2\sqrt{\bar\tau} \widetilde Y'(\bar\tau)\langle \widetilde
W(\bar\tau),X(\bar\tau)\rangle-\int_{\tau_1}^{\bar\tau}2\sqrt{\tau}\langle
\Jac(\widetilde Y'),W\rangle d\tau.$$ Since $W(\bar\tau)=0$ and
since $\nabla_{\widetilde Y'}(\widetilde W)=\nabla_W(\widetilde
Y')$, the boundary term vanishes. The integral vanishes since
$\widetilde Y'$ is an ${\mathcal L}$-Jacobi field.
\end{proof}

\subsection{Inequalities for $\triangle L_x^{\bar\tau}$}

The inequalities for the Hessian of $L_x^{\bar\tau}$ lead to
inequalities for $\triangle L_x^{\bar\tau}$ which we establish in this section. Here is the
main result.


\begin{proposition}[Laplacian inequality for $L_x^{\bar\tau}$]\label{prop:l-length-laplacian-inequality}
\uses{prop:l-length-hessian-inequality, def:l-injectivity-domain}
\label{triL}

Suppose that $q\in {\mathcal U}_x(\bar\tau)$, that $Z\in \widetilde
{\mathcal U}_x(\bar\tau)$ is the pre-image of $q$ and that
$\gamma_Z$ is the ${\mathcal L}$-geodesic determined by $Z$. Then
\begin{equation}\label{lineq}
 \triangle L_x^{\bar\tau}(q) \leq
\frac{n}{\sqrt{\bar\tau}-\sqrt{\tau_1}} -2\sqrt{\bar\tau}R(q) -
\frac{1}{(\sqrt{\bar\tau}-\sqrt{\tau_1})^2}{\mathcal
K}_{\tau_1}^{\bar\tau}(\gamma_Z),
\end{equation}
where, for any path $\gamma$ parameterized by backward time on the
interval
 $[\tau_1,\bar\tau]$ taking
 value $x$ at $\tau=\tau_1$ we define
$${\mathcal K}_{\tau_1}^{\bar\tau}(\gamma)= \int_{\tau_1}^{\bar\tau}
\sqrt{\tau}(\sqrt{\tau}-\sqrt{\tau_1})^2H(X)d\tau,$$ with
\begin{equation}\label{Hdefnn}
H(X)= -\frac{\partial R}{\partial \tau} - \frac{1}{\tau}R - 2\langle
\nabla R,X\rangle + 2\Ric(X,X), \end{equation} where $X$ is the
horizontal projection of $\gamma'(\tau)$. Furthermore,
Inequality~(\ref{lineq}) is an equality if and only if for every
$Y\in T_q(M_{T-\bar\tau})$ the unique adapted vector field $Y(\tau)$
along $\gamma$ satisfying $Y(\bar\tau)=Y$ is an ${\mathcal
L}$-Jacobi field. In this case
$$\Ric+\frac{1}{2\sqrt{\bar\tau}}\Hess(L_x^{\bar\tau})
=\frac{1}{2\sqrt{\bar\tau}(\sqrt{\bar\tau}-\sqrt{\tau_1})}G(T-\bar\tau).$$
\end{proposition}


\begin{proof}
  \uses{prop:l-length-hessian-inequality, lem:adapted-jacobi-hessian-relation, lem:div-ricci-scalar-curvature, thm:curvature-evolution}
\sketch

Choose an orthonormal basis $\{Y_{\alpha}\}$ for
$T_q(M_{T-\bar\tau})$. For each $\alpha$, extend $\{Y_{\alpha}\}$ to
an adapted vector field along the $ \mathcal{L}$-geodesic $\gamma_Z$
by solving
$$\nabla_{X}Y_{\alpha} = \frac{1}{2\sqrt{\tau}(\sqrt{\tau}-\sqrt{\tau_1})}Y_{\alpha}
-\Ric(Y_{\alpha},\cdot)^*.$$ As in Equation~(\ref{secondYeqn}),
we have
\begin{align*}
\frac{d}{d\tau}\langle Y_{\alpha},Y_{\beta}\rangle  &= \langle
\nabla_{X}Y_{\alpha},Y_{\beta}\rangle  + \langle
\nabla_{X}Y_{\beta}, Y_{\alpha}\rangle + 2\Ric(Y_{\alpha},Y_{\beta})
\\&= \frac{1}{\sqrt{\tau}(\sqrt{\tau}-\sqrt{\tau_1})}\langle Y_{\alpha},Y_{\beta}\rangle .
\end{align*}
 By integrating we get $$\langle Y_{\alpha},Y_{\beta}\rangle (\tau) =
\frac{(\sqrt{\tau}-\sqrt{\tau_1})^2}{(\sqrt{\bar\tau}-\sqrt{\tau_1})^2}\delta_{\alpha\beta}.$$
To simplify the notation we set
$$I(\tau)=\frac{\sqrt{\bar\tau}-\sqrt{\tau_1}}{\sqrt{\tau}-\sqrt{\tau_1}}$$
and $W_\alpha(\tau)=I(\tau)Y_\alpha(\tau)$.
 Then $\{W_{\alpha}(\tau)\}_\alpha$
form an orthonormal basis at $\tau$. Consequently, summing
Inequality~(\ref{Hessineq1}) over $\alpha$ gives
\begin{equation}\label{triLinequ}
\triangle L_x^{\bar\tau}(q)\le
\frac{n}{\sqrt{\bar\tau}-\sqrt{\tau_1}}
-2\sqrt{\bar\tau}R(q)-\sum_\alpha\int_{\tau_1}^{\bar\tau}\sqrt{\tau}H(X,Y_\alpha)d\tau.\end{equation}
To establish Inequality~(\ref{lineq}) it remains to prove the
following claim.


Clearly, Inequality~(\ref{lineq}) follows immediately from
Inequality~(\ref{triLinequ}) and the claim. The last statement of
Proposition~\ref{prop:l-length-laplacian-inequality} follows directly from the last statement of
Proposition~\ref{prop:l-length-hessian-inequality} and Lemma~\ref{lem:adapted-jacobi-hessian-relation}. This completes
the proof of Proposition~\ref{prop:l-length-laplacian-inequality}.
\end{proof}


\begin{claim}[Summed $H(X,Y_\alpha)$ formula]\label{claim:h-sum-formula}
\uses{prop:l-length-hessian-inequality}
\label{Hclaim}

$$\sum_\alpha
H(X,Y_\alpha)=\frac{(\sqrt{\tau}-\sqrt{\tau_1})^2}{(\sqrt{\bar\tau}-\sqrt{\tau_1})^2}H(X).$$
\end{claim}
\begin{proof}
\sketch
\uses{lem:div-ricci-scalar-curvature,thm:curvature-evolution}
To prove the claim we sum Equation~(\ref{Heqn}) giving
\begin{eqnarray*}
I^2(\tau) \sum_{\alpha}H(X,Y_{\alpha}) &= & \sum_\alpha H(X,W_\alpha)
\\
& = & -\triangle R + 2\Ric(X,X) - 4\langle \nabla R,X\rangle
+4\sum_{\alpha}(\nabla_{W_\alpha} \Ric)(W_{\alpha},X) \\ & &
-2\sum_{\alpha}\Ric_{\tau}( W_{\alpha},W_{\alpha}) +2|\Ric|^{2} -\frac{1}{\tau}R.
\end{eqnarray*}
Taking the trace of the second Bianchi identity gives
$$\sum_{\alpha}(\nabla_{W_\alpha}\Ric)(W_{\alpha},X)=\frac{1}{2}\langle\nabla R,X\rangle.$$ By~(\ref{Revol}),
$\partial R/\partial\tau=-\triangle R-2|\Ric|^2$, while differentiating
$R=g^{ij}R_{ij}$ gives $\sum_\alpha\partial_\tau\Ric(W_\alpha,W_\alpha)=-\triangle R$.
Thus $I^2(\tau)\sum_\alpha H(X,Y_\alpha)=H(X)$.
\end{proof}

\section{Reduced length}\label{sect:redln}

We introduce the reduced length function both on the tangent space
and on space-time. The reason that the reduced length $l_x$ is
easier to work with is that it is scale invariant when $\tau_1=0$.
Throughout Section~\ref{sect:redln} we fix $x\in{\mathcal M}$ with
${\bf t}(x)=T-\tau_1$. We shall always suppose that $T-\tau_1$ is
greater than the initial time of the generalized Ricci flow.

\subsection{The reduced length function $l_x$ on space-time}




\begin{definition}[Reduced length]\label{def:reduced-length}
\uses{def:minimal-l-length-function}

We define
 the $\mathcal{L}$-{\sl reduced length} (from $x$)
$$l_x\colon {\mathcal U}_x\to \Ar$$
by setting $$l_x(q)= \frac{L_x(q)}{2\sqrt{\tau}},$$ where $\tau=T-{\bf t}(q)$.
We denote by $l_x^\tau$ the restriction of $l_x$ to the slice ${\mathcal
U}_x(\tau)$.
\end{definition}


In order to understand the differential inequalities that $l_x$
satisfies, we first need to introduce a quantity closely related to
the function ${\mathcal K}^{\bar\tau}_{\tau_1}$ defined in
Proposition~\ref{prop:l-length-laplacian-inequality}.

\begin{definition}[The quantity $K_{\tau_1}^{\bar\tau}$]\label{def:reduced-length-k-quantity}
\uses{prop:l-length-laplacian-inequality}

For any ${\mathcal L}$-geodesic $\gamma$ parameterized by $[\tau_1,\bar\tau]$
we define
$$K_{\tau_1}^{\bar\tau}(\gamma)=
\int_{\tau_1}^{\bar\tau}\tau^{3/2}H(X)d\tau.$$
In the special case when $\tau_1=0$ we denote this integral by
$K^{\bar\tau}(\gamma)$.
\end{definition}

The following is immediate from the definitions.

\begin{lemma}[Continuity of the $K$-quantities]\label{lem:k-quantity-continuity}
\uses{def:reduced-length-k-quantity}
\label{Kequal}

For any ${\mathcal L}$-geodesic $\gamma$ defined on $[0,\bar\tau]$ both
$K^{\bar\tau}_{\tau_1}(\gamma)$ and ${\mathcal K}^{\bar\tau}_{\tau_1}(\gamma)$
are continuous in $\tau_1$ and at $\tau_1=0$ they take the same value. Also,
$$\left(\frac{\tau_1}{\bar\tau}\right)^{3/2}\left(R(\gamma(\tau_1))+|X(\tau_1)|^2\right)$$
is continuous for all $\tau_1> 0$ and has limit $0$ as
$\tau_1\rightarrow 0$. Here, as always, $X(\tau_1)$ is the
horizontal component of $\gamma'$ at $\tau=\tau_1$.
\end{lemma}

\begin{lemma}[Reduced length via $K$]\label{lem:reduced-length-k-formula}
\uses{def:reduced-length, def:reduced-length-k-quantity}
\label{Keqn}

Let $q\in {\mathcal
U}_x(\bar\tau)$, let $Z\in \widetilde {\mathcal U}_x$ be the
pre-image of $q$ and let $\gamma_Z$ be the ${\mathcal L}$-geodesic
determined by $Z$. Then we have
\begin{equation}\label{Kequat}
\bar\tau^{-\frac{3}{2}}K^{\bar\tau}_{\tau_1}(\gamma_Z)=\frac{
l_x(q)}{\bar\tau}-(R(q)
+|X(\bar\tau)|^2)+\left(\frac{\tau_1}{\bar\tau}\right)^{3/2}\left(R(x)+|X(\tau_1)|^2\right).
\end{equation}


In the case when $\tau_1=0$, the last term on the right-hand side of
Equation~(\ref{Kequat}) vanishes.
\end{lemma}


\begin{proof}
\sketch
\uses{prop:l-length-laplacian-inequality}
Using the ${\mathcal L}$-geodesic equation and the definition of $H$ we have
\begin{eqnarray*}
\lefteqn{\frac{d}{d\tau}(R(\gamma_Z(\tau))+ \abs{X(\tau)}^{2})} \\ &
= & \frac{\partial R}{\partial\tau}(\gamma_Z(\tau)) + \langle \nabla
R(\gamma_Z(\tau)),
X(\tau)\rangle + 2\langle \nabla_{X}X(\tau),X(\tau)\rangle \\
& & + 2\Ric(X(\tau),X(\tau))
\\&= & \frac{\partial R}{\partial\tau}(\gamma_Z(\tau)) + 2X(\tau)(R) - \frac{1}{\tau}\abs{X(\tau)}^{2} -
2\Ric(X(\tau),X(\tau)) \\&= &-H(X(\tau))
-\frac{1}{\tau}(R(\gamma_Z(\tau)+ \abs{X(\tau)}^{2}).
\end{eqnarray*}
Thus
$$\frac{d}{d\tau}(\tau^{\frac{3}{2}}(R(\gamma_Z(\tau)+\abs{X(\tau)}^{2}) )=
\frac{1}{2}\sqrt{\tau}(R(\gamma_Z(\tau)+\abs{X(\tau)}^{2}) -
\tau^{\frac{3}{2}}H(X(\tau)).$$ Integration from $\tau_1$ to $\bar\tau$ gives
$$\bar\tau^{3/2}\left(R(q))+|X(\bar\tau)|^2\right)-\tau_1^{3/2}(R(x)+|X(\tau_1)|^2)
=\frac{L_x^{\bar\tau}(q)}{2}-K_{\tau_1}^{\bar\tau}(\gamma_Z),$$
which is equivalent to Equation~(\ref{Kequat}). In the case when
$\tau_1=0$, the last term on the right-hand side vanishes since
$$\lim_{\tau\rightarrow 0}\tau^{3/2}|X(\tau)|^2=0.$$
\end{proof}



Now we come to the most general of the differential inequalities for
$l_x$ that will be so important in what follows. Whenever the
expression
$\left(\frac{\tau_1}{\bar\tau}\right)^{3/2}\left(R(x)+|X(\tau_1)|^2\right)$
appears in a formula, it is interpreted to be zero in the case when
$\tau_1=0$.


\begin{lemma}[General differential inequalities for $l_x$]\label{lem:reduced-length-differential-inequalities}
\uses{def:reduced-length, lem:reduced-length-k-formula, prop:l-length-laplacian-inequality}
\label{linequal}

For any  $q\in{\mathcal U}_x(\bar\tau)$, let $Z\in
\widetilde{\mathcal U}_x(\bar\tau)$ be the pre-image of $q$ and let
$\gamma_Z$ be the ${\mathcal L}$-geodesic determined by $Z$. Then we
have
\begin{eqnarray*} \frac{\partial l_x}{\partial \tau}(q) & = &
R(q)-\frac{l_x(q)}{\bar\tau}+\frac{K_{\tau_1}^{\bar\tau}(\gamma_Z)}{2\bar\tau^{3/2}}-\frac{1}{2}\left(\frac{\tau_1}{\bar\tau}
\right)^{3/2}\left(R(x)+|X(\tau_1)|^2\right) \\
|\nabla l^{\bar\tau}_x(q)|^2 & = & |X(\bar\tau)|^2
=\frac{l^{\bar\tau}_x(q)}{\bar\tau}-\frac{K_{\tau_1}^{\bar\tau}
(\gamma_Z)}{\bar\tau^{3/2}}-R(q) +\left(\frac{\tau_1}{\bar\tau}
\right)^{3/2}\left(R(x)+|X(\tau_1)|^2\right) \\
\triangle l^{\bar\tau}_x(q) & = & \frac{1}{2\sqrt{\bar\tau}}\triangle
L^{\bar\tau}_x(q)\le \frac{n}{2\sqrt{\bar\tau}(\sqrt{\bar\tau}-\sqrt{\tau_1})}
-R(q)-\frac{{\mathcal
K}_{\tau_1}^{\bar\tau}(\gamma_Z)}{2\sqrt{\bar\tau}(\sqrt{\bar\tau}-\sqrt{\tau_1})^{2}}.
\end{eqnarray*}
\end{lemma}

\begin{proof}
\sketch
\uses{lem:l-length-tau-derivative,cor:l-length-gradient,lem:reduced-length-k-formula,prop:l-length-laplacian-inequality}
 It follows immediately from Equation~(\ref{Ltaueqn}) that
$$\frac{\partial l_x}{\partial \tau}=R-\frac{1}{2}(R+|X|^2)-\frac{l_x}{2\tau}.$$
Using Equation~(\ref{Kequat}) this gives the first equality stated
in the lemma.
 It follows immediately from Corollary~\ref{cor:l-length-gradient} that
$\nabla l_x^\tau=X(\tau)$ and hence $|\nabla
l_x^\tau|^2=|X(\tau)|^2$. From this and Equation~(\ref{Kequat}) the
second equation follows. The last inequality is immediate from
Proposition~\ref{prop:l-length-laplacian-inequality}.

When $\tau_1=0$, the last terms on the right-hand sides of the first
two equations vanish, since the last term on the right-hand side of
Equation~(\ref{Kequat}) vanishes in this case.
\end{proof}

When $\tau_1=0$, which is the case of main interest, all these
formulas simplify and we get:

\begin{theorem}[Reduced length differential identities at $\tau_1=0$]\label{thm:reduced-length-differential-identities}
\uses{def:reduced-length, lem:reduced-length-differential-inequalities}
\label{Ueqn}

Suppose that $x\in M_T$ so that $\tau_1=0$. For any $q\in{\mathcal
U}_x(\bar\tau)$, let $Z\in \widetilde{\mathcal U}_x(\bar\tau)$ be
the pre-image of $q$ and let $\gamma_Z$ be the ${\mathcal
L}$-geodesic determined by $Z$. As usual, let $X(\tau)$ be the
horizontal projection of $\gamma_Z'(\tau)$. Then we have
\begin{eqnarray*} \frac{\partial l_x}{\partial \tau}(q) & = &
R(q)-\frac{l_x(q)}{\bar\tau}+\frac{K^{\bar\tau}(\gamma_Z)}{2\bar\tau^{3/2}}
\\
|\nabla l_x^{\bar\tau}(q)|^2 & = & |X(\bar\tau)|^2  =
\frac{l_x^{\bar\tau}(q)}{\bar\tau}-\frac{K^{\bar\tau}
(\gamma_Z)}{\bar\tau^{3/2}}-R(q)  \\
\triangle l_x^{\bar\tau}(q) & = &
\frac{1}{2\sqrt{\bar\tau}}\triangle L_x^{\bar\tau}(q)\le \frac{n}{2
\bar\tau}-R(q)-\frac{K^{\bar\tau}(\gamma_Z)}{2\bar\tau^{3/2}}.
\end{eqnarray*}
\end{theorem}

\begin{proof}
\sketch
\uses{lem:reduced-length-differential-inequalities}
This is immediate from the formulas in the previous lemma.
\end{proof}

Now let us reformulate the differential inequalities in
Theorem~\ref{thm:reduced-length-differential-identities} in a way that will be useful later.

\begin{corollary}[Reformulated PDE inequalities for $l_x$]\label{cor:reduced-length-pde-inequalities}
\uses{thm:reduced-length-differential-identities}
\label{lformula}

Suppose that $x\in M_T$ so that $\tau_1=0$. Then for $q\in{\mathcal
U}^{\bar\tau}_x$ we have
$$\frac{\partial l_x}{\partial \tau}(q)+\triangle l_x^{\bar\tau}(q)\le \frac{(n/2)-l_x^{\bar\tau}(q)}{\bar\tau}.$$
$$\frac{\partial l_x}{\partial \tau}(q)-\triangle l_x^{\bar\tau}(q)+|\nabla l_x^{\bar\tau}(q)|^2-R(q)+\frac{n}{2\bar\tau}\ge 0 .$$
$$2\triangle l_x^{\bar\tau}(q)-|\nabla l_x^{\bar\tau}(q)|^2+R(q)+\frac{l_x^{\bar\tau}(q)-n}{\bar\tau}
\le 0.$$

In fact, setting $$\delta=\frac{n}{2
\bar\tau}-R(q)-\frac{K^{\bar\tau}(\gamma_Z)}{2\bar\tau^{3/2}}-\triangle
l_x^{\bar\tau}(q),$$ then $\delta\ge 0$ and
$$\frac{\partial l_x}{\partial \tau}(q)-\triangle l_x^{\bar\tau}(q)
+|\nabla l_x^{\bar\tau}(q)|^2-R(q)+\frac{n}{2\bar\tau}=\delta$$
$$2\triangle l_x^{\bar\tau}(q)-|\nabla l_x^{\bar\tau}(q)|^2+R(q)
+\frac{l_x^{\bar\tau}(q)-n}{\bar\tau}=-2\delta.$$
\end{corollary}

\subsection{The tangential version $\widetilde l$ of the reduced length function}

For any path $\gamma\colon [\tau_1,\bar\tau]\to ({\mathcal M},G)$
parameterized by backward time we define
$$l(\gamma)=\frac{1}{2\sqrt{\bar\tau}}{\mathcal L}(\gamma).$$
This leads immediately to a reduced length on $\widetilde {\mathcal
U}_x$.


\begin{definition}[Tangential reduced length $\widetilde l$]\label{def:reduced-length-tangential}
\uses{def:l-length, def:l-exponential-map}

We define $\widetilde l\colon \widetilde {\mathcal U}_x\to
\Ar$  by
$$\widetilde l(Z,\tau)=\frac{\widetilde L(Z,\tau)}{2\sqrt{\tau}}=l(\gamma_Z|_{[\tau_1,\bar\tau]}).$$
\end{definition}

At first glance it may appear that the computations of the gradient
and $\tau$-derivatives for $l_x$ pass immediately to those for
$\widetilde l$. For the spatial derivative this is correct, but for
the $\tau$-derivative it is not true. As the computation below
shows, the $\tau$-derivatives of $\widetilde l$ and $l_x$ do not
agree under the identification ${\mathcal L}\exp_x$. The reason
is that this identification does not line up the $\tau$-vector field
in the domain with $-\partial/\partial t$ in the range. So it is an
entirely different computation with a different answer.




\begin{lemma}[$\tau$-derivative of $\widetilde l$]\label{lem:reduced-length-tangential-tau-derivative}
\uses{def:reduced-length-tangential}
\label{tildeltau}

$$\frac{\partial \widetilde
l(Z,\tau)}{\partial
\tau}=\frac{1}{2}\left(R(\gamma_Z(\tau))+X(\tau)|^2\right)-\frac{\widetilde
l(Z,\tau)}{2\tau}.$$
\end{lemma}

\begin{proof}
\sketch
\uses{def:reduced-length-tangential}
By the Fundamental Theorem of Calculus
$$\frac{\partial}{\partial\tau}\widetilde L(Z,\tau)=\sqrt{\tau}\left(R(\gamma_z(\tau))+|X(\tau)|^2\right).$$
Thus,
$$\frac{\partial}{\partial\tau}\widetilde l(Z,\tau)=\frac{1}{2}\left(R(\gamma_z(\tau))+|X(\tau)|^2\right)
-\frac{\widetilde l(Z,\tau)}{2\tau}.$$
\end{proof}

\begin{corollary}[$\tau$-derivative of $\widetilde l$ via $K$]\label{cor:reduced-length-tangential-k-derivative}
\uses{lem:reduced-length-tangential-tau-derivative, lem:reduced-length-k-formula}
\label{tildelK}

Suppose that $x\in M_T$ so that
$\tau_1=0$. Then
$$\frac{\partial}{\partial\tau}\widetilde
l(Z,\tau)=-\frac{K^\tau(\gamma_Z)}{2\tau^{\frac{3}{2}}}.$$
\end{corollary}

\begin{proof}
\sketch
\uses{lem:reduced-length-tangential-tau-derivative,lem:reduced-length-k-formula}
This is immediate from Lemma~\ref{lem:reduced-length-tangential-tau-derivative} and Lemma~\ref{lem:reduced-length-k-formula} (after the
latter is rewritten using $\widetilde L$ instead of $L_x$).
\end{proof}



\section{Local Lipschitz estimates for $l_x$}\label{sect:lips}

It is important for the applications to have results on the
Lipschitz properties of $l_x$, or equivalently $L_x$. Of course,
these are the analogues of the fact that in Riemannian geometry the
distance function from a point is Lipschitz. The proof of the
Lipschitz property given here is based on the exposition in
\cite{Ye}. In Section~\ref{sect:lips}, we fix $x\in
M_{T-\tau_1}\subset {\mathcal M}$.

\subsection{Statement and corollaries}


\begin{definition}[Domain of definition of $l_x$]\label{def:reduced-length-domain}
\uses{def:reduced-length}

Let $({\mathcal M},G)$ be a generalized Ricci flow and let $x\in
M_{T-\tau_1}\subset {\mathcal M}$. The reduced length function $l_x$ is defined
on the subset of ${\mathcal M}$ consisting of all points $y\in {\mathcal M}$
for which there is a minimizing ${\mathcal L}$-geodesic from $x$ to $y$. The
value $l_x(y)$ is the quotient of ${\mathcal L}$-length of any such minimizing
${\mathcal L}$-geodesic divided by $2\sqrt{\tau}$.
\end{definition}

Here is the main result of this subsection.

\begin{proposition}[Local Lipschitz property of $l_x$]\label{prop:reduced-length-lipschitz}
\uses{def:reduced-length-domain}
\label{lips}

Let $({\mathcal M},G)$ be a generalized Ricci flow and let $x\in
M_{T-\tau_1}\subset {\mathcal M}$. Let $\epsilon>0$ be given and let $A\subset
{\mathcal M}\cap {\bf t}^{-1}(-\infty,T-\tau_1+\epsilon)$. Suppose that there
is a subset $F\subset {\mathcal M}$ on which $|\Ric|$ and $|\nabla R|$ are
bounded and a neighborhood $\nu(A)$ of $A$ contained in $F$ with the property
that for every point $z\in \nu(A)$ there is a minimizing ${\mathcal
L}$-geodesic from $x$ to $z$ contained in $F$. Then $l_x$ is defined on  all of
$\nu(A)$. Furthermore, there is a smaller neighborhood $\nu_0(A)\subset \nu(A)$
of $A$ on which $l_x$ is a locally Lipschitz function with respect to the
Riemannian metric, denoted $\widehat G$, on ${\mathcal M}$ which is defined as
the orthogonal sum of the Riemannian metric $G$ on ${\mathcal H}T{\mathcal M}$
and the metric $dt^2$ on the tangent line spanned by $\chi$.
\end{proposition}







\begin{corollary}[Lipschitz property on a time-slice]\label{cor:reduced-length-lipschitz-time-slice}
\uses{prop:reduced-length-lipschitz}

With $A$ and $\nu_0(A)$ as in Proposition~\ref{prop:reduced-length-lipschitz}, the
restriction of $l_x$ to  $\nu_0(A)\cap M_{T-\bar\tau}$ is a locally
Lipschitz function with respect to the metric $G_{T-\bar\tau}$.
\end{corollary}


\subsection{The proof of Proposition~{\ref{prop:reduced-length-lipschitz}}}

Proposition~\ref{prop:reduced-length-lipschitz} follows from a much more precise, though more
complicated to state, result. In order to state this more technical
result we introduce the following definition.

\begin{definition}[Parabolic-type neighborhood $\widetilde P$]\label{def:l-parabolic-neighborhood}
\uses{def:compatible-embedding}

Let $y\in {\mathcal M}$ with ${\bf t}(y)=t$ and suppose that for
some $\epsilon>0$ there is an embedding $\iota\colon B(y,t,r)\times
(t-\epsilon,t+\epsilon)\to {\mathcal M}$ that is compatible with
time and the vector field. Then we denote by $\widetilde
P(y,r,\epsilon)\subset {\mathcal M}$ the image of $\iota$. Whenever
we introduce $\widetilde P(y,r,\epsilon)\subset {\mathcal M}$
implicitly we are asserting that such an embedding exists.

For  $A\subset {\mathcal M}$, if $\widetilde
P(a,\epsilon,\epsilon)\subset {\mathcal M}$ exists for every $a\in
A$, then we denote by $\nu_\epsilon(A)$ the union over all $a\in A$
of $\widetilde P(a,\epsilon,\epsilon)$.
\end{definition}


Now we are ready for the more precise, technical result.


\begin{proposition}[Precise local Lipschitz estimate for $l_x$]\label{prop:reduced-length-lipschitz-precise}
\uses{def:l-parabolic-neighborhood, def:reduced-length-domain}
\label{lipaty}

Given constants  $\epsilon>0$, $\bar\tau_0<\infty$, $l_0<\infty$,
and $C_0<\infty$, there are constants $ C<\infty$ and
$0<\delta<\epsilon$ depending only the  given constants  such that
the following holds. Let $({\mathcal M},G)$ be a generalized Ricci
flow and let $x\in{\mathcal M}$ be a point with ${\bf
t}(x)=T-\tau_1$. Let $y\in {\mathcal M}$ be a point with ${\bf
t}(y)=t=T-\bar\tau$ where $\tau_1+\epsilon\le \bar\tau\le
\bar\tau_0$. Suppose that there is a minimizing ${\mathcal
L}$-geodesic $\gamma$ from $x$ to $y$ with $l(\gamma)\le l_0$.
Suppose that the ball $B(y,t,\epsilon)$ has compact closure in
$M_{t}$ and that $\widetilde P(y,\epsilon,\epsilon)\subset {\mathcal
M}$ exists and that the sectional curvatures of the restriction of
$G$ to this submanifold are bounded by $C_0$. Lastly, suppose that
for every point of the form $z\in \widetilde P(y,\delta,\delta)$
there is a minimizing ${\mathcal L}$-geodesic from $x$ to $z$ with
$|\Ric|$ and $|\nabla R|$ bounded by $C_0$ along this geodesic.
Then for all $(b,t')\in B(y,t,\delta)\times(t-\delta,t+\delta)$ we
have
$$|l_x(y)-l_x(\iota(b,t'))|\le C\sqrt{d_{t}(y,b)^2+|t-t'|^2}.$$
\end{proposition}




Before proving Proposition~\ref{prop:reduced-length-lipschitz-precise}, let us show how it implies
Proposition~\ref{prop:reduced-length-lipschitz}.

\begin{proof}
\sketch
Proposition~\ref{prop:reduced-length-lipschitz})
\uses{prop:reduced-length-lipschitz-precise,def:l-parabolic-neighborhood}
Suppose given $\epsilon>0$, $A$, $\nu(A)$ and $F$ as in
the statement of Proposition~\ref{prop:reduced-length-lipschitz}. For each $y\in A$ there is
$0<\epsilon'<\epsilon$ and a neighborhood $\nu'(y)$ with (i) the closure
$\overline{\nu}'(y)$ of $\nu'(y)$ being a compact subset of $\nu(A)$ and (ii)
for each $z\in \overline{\nu}'(y)$ the parabolic neighborhood $\widetilde
P(z,\epsilon',\epsilon')$ exists and has compact closure in $\nu(A)$. It
follows that for every $z\in \overline{\nu}'(A)$, $\Rm_G$ is bounded on
$\widetilde P(z,\epsilon',\epsilon')$ and every point of $\widetilde
P(z,\epsilon',\epsilon')$ is connected to $x$ by a minimizing ${\mathcal
L}$-geodesic in $F$. Thus, Proposition~\ref{prop:reduced-length-lipschitz-precise}, with $\epsilon$ replaced
by $\epsilon'$, applies to $z$. In particular, $l_x$ is continuous at $z$, and
hence is continuous on all of $\overline{\nu}'(y)$. Thus, $l_x$ is bounded on
$\overline{\nu}'(y)$. Since we have uniform bounds for the curvature on
$\widetilde P(z,\epsilon',\epsilon')$ according to Proposition~\ref{prop:reduced-length-lipschitz-precise}
there are constants $C<\infty$ and $0<\delta<\epsilon'$ such that for any $z\in
\overline{\nu}'(y)$ and any $z'\in \widetilde P(z,\delta,\delta)$, we have
$$|l_x(z)-l_x(z')|\le C|z-z'|_{G({\bf t}(z)+dt^2}.$$
Since we have a uniform bound for the curvature on $\widetilde
P(z,\epsilon',\epsilon')$ independent of $z\in \nu'(y)$, the metrics $\widehat
G=G+dt^2$ and $G({\bf t}(z))+dt^2$ are uniformly comparable on all of
$\widetilde P(z,\delta,\delta)$. It follows that there is a constant
$C'<\infty$ such that for all $z\in \nu'(y)$ and all $z'\in \widetilde
P(z,\delta,\delta)$ we have
$$|l_x(z)-l_x(z')|\le C'|Z-z'|_{\widehat G}.$$

We set $\nu_0(A)=\cup_{y\in A}\nu'(y)$. This is an open neighborhood of $A$
contained in $\nu(A)$ on which $l_x$ is locally Lipschitz with respect to the
metric $\widehat G$.
\end{proof}



Now we turn to the proof of Proposition~\ref{prop:reduced-length-lipschitz-precise}. We begin with
several preliminary results.



\begin{lemma}[Min-max estimate along an ${\mathcal L}$-geodesic]\label{lem:l-geodesic-min-max-estimate}
\uses{def:l-geodesic}
\label{min-max}

 Suppose  that
$\gamma$ is an ${\mathcal L}$-geodesic defined on
$[\tau_1,\bar\tau]$, and suppose that for all $\tau\in
[\tau_1,\bar\tau]$ we have $|\nabla R(\gamma(\tau))|\le C_0$ and
$|\Ric(\gamma(\tau))|\le C_0$. Then
$$\max_{\tau}\left(\sqrt{\tau}|X_\gamma(\tau)|\right)\le C_1\min_{\tau}\left(\sqrt{\tau}|X_\gamma(\tau)|\right)+\frac{(C_1-1)}{2}\sqrt{\bar\tau},$$
where $C_1=e^{2C_0\bar\tau}$.
\end{lemma}

\begin{proof}
\sketch
\uses{def:l-geodesic}
The geodesic equation in terms of the variable $s$,
Equation~(\ref{reparam}), gives
\begin{eqnarray}
\frac{d|\gamma'(s)|^2}{ds} & = & 2\langle
\nabla_{\gamma'(s)}\gamma'(s),\gamma'(s)\rangle+4s\Ric(\gamma'(s),\gamma'(s))\nonumber \\
& = & 4s^2\langle\nabla R,\gamma'(s)\rangle -4s\Ric(\gamma'(s),\gamma'(s))\label{nablas}.
\end{eqnarray}

Thus,, by our assumption on $|\nabla R|$ and $|\Ric|$ along
$\gamma$, we have
$$\left|\frac{d|\gamma'(s)|^2}{ds}\right| \le
4C_0s^2|\gamma'(s)|+4C_0s|\gamma'(s)|^2.$$ It follows that
$$\left|\frac{d|\gamma'(s)|}{ds}\right|\le 2C_0s^2+2C_0s|\gamma'(s)|\le 2C_0\bar\tau+2
C_0\sqrt{\bar\tau}|\gamma'(s)|,$$ and hence that
$$-2C_0\sqrt{\bar\tau}ds\le \frac{d|\gamma'(s)|}{\sqrt{\bar\tau}+|\gamma'(s)|}
\le 2C_0\sqrt{\bar\tau}ds.$$ Suppose that $s_0<s_1$. Integrating from $s_0$ to
$s_1$ gives
\begin{eqnarray*}
|\gamma'(s_1)| & \le &  C|\gamma'(s_0)|+(C-1)\sqrt{\bar\tau} \\
|\gamma'(s_0)| & \le & C|\gamma'(s_1)|+(C-1)\sqrt{\bar\tau}
\end{eqnarray*}
where
\begin{equation*}
C=e^{2C_0\sqrt{\bar\tau}(s_1-s_0)}.
\end{equation*}
 Since
$\sqrt{\tau}X_\gamma(\tau)=\frac{1}{2}\gamma'(s)$, this completes
the proof of the lemma.
\end{proof}

\begin{corollary}[Uniform speed bound along an ${\mathcal L}$-geodesic]\label{cor:l-geodesic-speed-bound}
\uses{lem:l-geodesic-min-max-estimate, def:l-length}
\label{minmax}

Given $\bar\tau_0<\infty$, $C_0<\infty$, $\epsilon>0$, and $l_0<\infty$, there
is a constant $C_2$ depending only on $C_0$, $l_0$, $\epsilon$ and $\bar\tau_0$
such that the following holds. Let $\gamma$ be an ${\mathcal L}$-geodesic
defined on $[\tau_1,\bar\tau]$ with $\tau_1+\epsilon\le\bar\tau\le \bar\tau_0$
and with $|\nabla R(\gamma(\tau))|\le C_0$ and $|\Ric(\gamma(\tau))|\le
C_0$ for all $\tau\in [\tau_1,\bar\tau]$. Suppose also that $ l(\gamma)\le
l_0$. Then, we have
$$\max_{\tau}\left(\sqrt{\tau}|X_\gamma(\tau)|\right)\le C_2.$$
\end{corollary}

\begin{proof}
\sketch
\uses{lem:l-geodesic-min-max-estimate}
From the definition
${\mathcal L}(\gamma)=\int_{\sqrt{\tau_1}}^{\sqrt{\bar\tau}}
(2s^2R+\frac{1}{2}|\gamma'(s)|^2)ds$. Because of the bound on $|\Ric|$ (which implies that $|R|\le 3C_0$) we have
$$\frac{1}{2}\int_{\sqrt{\tau_1}}^{\sqrt{\bar\tau}}|\gamma'(s)|^2ds\le {\mathcal L}(\gamma)+
2C_0\bar\tau^{3/2}.$$ Thus,
$$(\sqrt{\bar \tau}-\sqrt{\tau_1})\min(|\gamma'(s)|^2)\le 2{\mathcal L}(\gamma)+
4C_0\bar\tau^{3/2}.$$ The bounds $\tau_1+\epsilon\le \bar\tau\le\bar\tau_0$,
then imply that $\min|\gamma'(s)|^2\le C''$ for some $C''$ depending on
$C_0,l_0,\epsilon,$ and $\bar\tau_0$. Since
$\sqrt{\tau}X_\gamma(\tau)=\frac{1}{2}\gamma'(s)$, we have
$$\min_{\tau}\left(\sqrt{\tau}|X_\gamma(\tau)|\right)\le C'$$ for some constant $C'$ depending
only on $C_0$, $l_0$, $\epsilon$ and $\bar\tau_0$. The result is now immediate
from Lemma~\ref{lem:l-geodesic-min-max-estimate}.
\end{proof}

Now we are ready to show that, for $z$ sufficiently close to  $y$,
the reduced length $l_x(z)$ is bounded above by a constant depending
on the curvature bounds, on $l_x(y)$, and on the distance in
space-time from $z$ to $y$.


\begin{lemma}[Upper bound for $l_x$ near a given point]\label{lem:reduced-length-upper-bound-nearby}
\uses{cor:l-geodesic-speed-bound, def:l-parabolic-neighborhood, def:l-length}
\label{upperbd}

Given constants  $\epsilon>0$, $\bar\tau_0<\infty$, $C_0<\infty$,
and $l_0<\infty$, there are $C_3<\infty$ and $0<\delta_2\le
\epsilon/4$ depending only on the  given constants  such that the
following holds. Let $y\in {\mathcal M}$ be a point with ${\bf
t}(y)=t_0=T-\bar\tau$ where $\tau_1+\epsilon\le \bar\tau\le
\bar\tau_0$. Suppose that there is a minimizing ${\mathcal
L}$-geodesic $\gamma$ from $x$ to $y$ with $l_x(\gamma)\le l_0$.
Suppose that $|\nabla R|$ and $|\Ric|$ are bounded by $C_0$
along $\gamma$. Suppose also that the ball $B(y,t_0,\epsilon)$ has
compact closure in $M_{t_0}$ and that there is an embedding
$$\iota\colon B(y,t_0,\epsilon)\times (t_0-\epsilon,t_0+\epsilon)\stackrel
{\cong}{\longrightarrow} \tilde P(y,\epsilon,\epsilon)\subset
{\mathcal M}$$ compatible with  time  and the vector field so that
the sectional curvatures of the restriction of $G$ to the image of
this embedding are bounded by $C_0$. Then for any point $b\in
B(y,t_0,\delta_2)$ and for any $t'\in (t_0-\delta_2,t_0+\delta_2)$
there is a curve $\gamma_1$ from $x$ to the point $z=\iota(b,t')$,
parameterized by backward time, such that
$$l(\gamma_1)\le  l(\gamma)+C_3\sqrt{d_{t_0}(y,b)^2+|t_0-t'|^2}.$$
\end{lemma}

\begin{proof}
\sketch
\uses{cor:l-geodesic-speed-bound}
Let $C_2$ be the constant depending on $C_0$, $l_0$, $\epsilon$, and
$\bar\tau_0$ from Corollary~\ref{cor:l-geodesic-speed-bound}, and set
$$C'=\frac{\sqrt{2}}{\sqrt{\epsilon}}C_2.$$
Since $\bar\tau\ge \epsilon$, it follows that
$\bar\tau-\epsilon/2\ge \epsilon/2$, so that by
Corollary~\ref{cor:l-geodesic-speed-bound} we have $|X_\gamma(\tau)|\le C'$ for all
$\tau\in [\bar\tau-\epsilon/2,\bar\tau]$. Set $0<\delta_0$
sufficiently small (how small depends only on $C_0$) such that for
all $(z,t)\in \tilde P(y,\epsilon,\delta_0)$ we have
$$\frac{1}{2}g(z,t)\le g(z,t_0)\le 2g(z,t),$$ and define
$$\delta_2=\min\left(\frac{\epsilon}{8},\frac{\epsilon}{8C'},\frac{\delta_0}{4}\right).$$
 Let $b\in B(y,t_0,\delta_2)$ and $t'\in
(t_0-\delta_2,t_0+\delta_2)$ be given. Set
$\alpha=\sqrt{d_{t_0}(y,b)^2+|t_0-t'|^2}$, set $t_1=t_0-2\alpha$,
and let $\tau_1=T-t_1$. Notice that
$\alpha<\sqrt{2}\delta_2<\epsilon/4$, so that the norm of the Ricci
curvature is bounded by $C'$ on
$\iota(B(y,t_0,\epsilon)\times(t_1,t_0+2\alpha))$.



Now let $\overline\mu\colon [\bar\tau-2\alpha,T-t']\to
B(y,t_0,\epsilon)$ be a shortest $g(t_0)$-geodesic from $c$ to $b$,
parameterized at constant $g(t_0)$-speed, and let $\mu$ be the path
parameterized by backward time defined by
$$\mu(\tau)=\iota(\overline\mu(\tau),T-\tau)$$
for all $\tau\in [\bar\tau-2\alpha,T-t']$. Then the concatenation
$\gamma_1=\gamma|_{[\tau_1,\bar\tau-2\alpha]}*\mu$ is a path
parameterized by backward time from $x$ to $\iota(b,t')$.





On the other hand, since $R\ge -3C_0$ in $P(y,\epsilon,\epsilon)$
and $|X|^2\ge 0$, we see that
$${\mathcal L}(\gamma|_{[\tau_1,\bar\tau-2\alpha]})\le {\mathcal L}(\gamma)
+ 6C_0\alpha\sqrt{\bar\tau_0}.$$  Together with the previous claim
this establishes Lemma~\ref{lem:reduced-length-upper-bound-nearby}.
\end{proof}

\begin{claim}[${\mathcal L}$-geodesic endpoint stays close to the ball]\label{claim:l-geodesic-endpoint-close-to-ball}
\uses{cor:l-geodesic-speed-bound}
\label{6.66}

$\gamma(\tau_1)\in \widetilde P(y,\epsilon,\epsilon)$ and writing
$\gamma(\tau_1)=\iota(c,t_1)$ we have $d_{t_0}(c,b)\le
(4C'+1)\alpha$.
\end{claim}

\begin{proof}
\sketch
Since $|X_\gamma(\tau)|\le C'$ for all $\tau\in
[\bar\tau-2\alpha,\bar\tau]$, and  $\delta_2\le \delta_0/4$, it
follows that $2\alpha\le \delta_0$ and hence that
$|X_\gamma(\tau)|_{g(t_0)}\le 2C'$ for all  $\tau\in
[\bar\tau-2\alpha,\bar\tau]$. Since $\gamma(\bar\tau)=y$, this
implies that
$$d_{t_0}(y,c)\le 4C'\alpha.$$
The claim then follows from the triangle inequality.
\end{proof}

\begin{claim}[${\mathcal L}$-length of the concatenated path]\label{claim:l-length-concatenation-bound}
\uses{claim:l-geodesic-endpoint-close-to-ball}

There is a constant $C'_1$ depending only on
$C_0$, $C'$, and $\bar\tau_0$ such that
$$l(\gamma_1)\le l(\gamma|_{[\tau_1,\bar\tau-2\alpha]})+C'_1\alpha$$
\end{claim}

\begin{proof}
  \uses{lem:reduced-length-upper-bound-nearby}
\sketch
First notice that since $\bar\tau=T-t_0$ and $|t'-t_0|\le \alpha$ we have
$(T-t')-(\bar\tau-2\alpha)=2\alpha+(t'-t_0)\ge \alpha$. According to
Claim~\ref{claim:l-geodesic-endpoint-close-to-ball} this implies that the $g(t_0)$-speed of $\mu$ is at most
$(4C'+1)$, and hence  that $|X_{\mu}(\tau)|_{g(T-\tau)}\le 8C'+2$ for all
$\tau\in [\bar\tau-2\alpha,T-t']$. Consequently, $R+|X_\mu|^2$ is bounded above
along $\mu$ by a constant $\widetilde C$ depending only on $C'$ and $C_0$. This
implies that ${\mathcal L}(\mu)\le \widetilde C\alpha\sqrt{T-t'}$. Of course,
$T-t'\le \bar\tau+\epsilon<2\bar\tau\le 2\bar\tau_0$. This completes the proof
of the claim.
\end{proof}

This is a one-sided inequality which says that the nearby values of
$l_x$ are bounded above in terms of  $l_x(y)$, the curvature bounds,
and the distance in space-time from $y$. In order to complete the
proof of Proposition~\ref{prop:reduced-length-lipschitz-precise} we must establish inequalities in
the opposite direction. This requires reversing the roles of the
points.


\begin{proof}
\sketch
\uses{lem:reduced-length-upper-bound-nearby}
Let $\delta_2$ and $C_3$ be the constants given by
Lemma~\ref{lem:reduced-length-upper-bound-nearby} associated to $\epsilon/2$, $\bar\tau_0$, $C_0$,
and $l_0$. We shall choose $C\ge C_3$ and $\delta\le \delta_2$ so
that by Lemma~\ref{lem:reduced-length-upper-bound-nearby} we will automatically have
$$l_x(\iota(b,t'))\le l_x(y)+C_3\sqrt{d_{t_0}(y,b)^2+|t_0-t'|^2}
\le l_x(y)+C\sqrt{d_{t_0}(y,b)^2+|t_0-t'|^2}$$ for all
$\iota(b,t')\in P(y,\delta,\delta)$. It remains to choose $C$ and
$\delta$ so that
$$l_x(y)\le l_x(\iota(b,t'))+C\sqrt{d_{t_0}(y,b)^2+|t_0-t'|^2}.$$
 Let
$\delta'_2$ and $ C'_3$ be the constants given by
Lemma~\ref{lem:reduced-length-upper-bound-nearby} for the following set of input constants:
$C_0'=C_0$, $\bar\tau_0$ replaced by
$\bar\tau_0'=\bar\tau_0+\epsilon/2$, and $l_0$ replaced by
$l_0'=l_0+\sqrt{2}C_3\delta_2$, and $\epsilon$ replaced by
$\epsilon'=\epsilon/4$. Then set $C=\max(2C_3',C_3)$.

Let $z=\iota(b,t')\in \widetilde P(y,\delta,\delta)$.



 By
the above and the fact that $ \delta\le\epsilon/4$, the sectional
curvatures on $\widetilde P(z,\epsilon/4,\epsilon/4)$ are bounded by
$C_0$. By Lemma~\ref{lem:reduced-length-upper-bound-nearby} there is a curve parameterized by
backward time from $x$ to $z$ whose $l$-length is at most $l_0'$.
Thus the $l$-length any minimizing ${\mathcal L}$-geodesic from $x$
to $z$ is at most $l_0'$. By assumption we have a minimizing
${\mathcal L}$-geodesic with the property that  $|\Ric|$
and$|\nabla R|$ are bounded by $C_0$ along the ${\mathcal
L}$-geodesic.

Of course, $t_0-\delta<t'<t_0+\delta$, so that
$\tau_1+\epsilon/2<T-t'\le \bar\tau_0+\epsilon/4$. This means that
Lemma~\ref{lem:reduced-length-upper-bound-nearby} applies to show that for very $w=\iota(c,t)\in
\widetilde P(z,\delta_2',\delta'_2)$,  we have
$$l_x(w)\le l_x(z)+C'_3\sqrt{d_{t'}(b,c)^2+|t-t'|^2}.$$


The proof is then completed by showing the following:



It follows immediately that \begin{eqnarray*} l_x(y) & \le &
l_x(z)+C'_3\sqrt{d_{t'}(b,y)^2+|t_0-t'|^2} \\ & \le & l_x(z)+2
C'_3\sqrt{d_{t_0}(b,y)^2+|t_0-t'|^2}\le
l_x(z)+C\sqrt{d_{t_0}(b,y)^2+|t_0-t'|^2}.
\end{eqnarray*}

This completes the proof of Proposition~\ref{prop:reduced-length-lipschitz-precise}.
\end{proof}

\begin{claim}[The parabolic ball stays inside $B(y,t_0,\epsilon)$]\label{claim:parabolic-ball-containment}
For $\delta$ sufficiently small (how small depending on $\delta_2$
and $\delta'_2$) we have $B(z,t',\epsilon/4)\subset
B(y,t_0,\epsilon)$.
\end{claim}

\begin{proof}
\sketch
Since $|t_0-t'|<\delta\le \delta_2$, and by construction
$\delta_2<\delta_0$, it follows  that for any $c\in
B(y,t_0,\epsilon)$ we have $d_{t'}(b,c)\le 2d_{t_0}(b,c)$. Since
$d_{t_0}(y,b)<\delta\le \epsilon/4$, the result is immediate from
the triangle inequality.
\end{proof}

\begin{claim}[$y$ lies in the parabolic neighborhood of $z$]\label{claim:point-in-parabolic-neighborhood}
$y\in \widetilde P(z,\delta_2',\delta'_2)$.
\end{claim}

\begin{proof}
\sketch
By construction $|t'-t_0|<\delta\le \delta_2'$. Also,
$d_{t_0}(y,b)<\delta\le \delta'_2/2$. Since $d_{t_0}\le 2d_{t'}$, we
have $d_{t'}(y,b)< \delta'_2$ the claim is then immediate.
\end{proof}

\begin{corollary}[Full measure of the injectivity set]\label{cor:reduced-length-full-measure}
\uses{def:l-injectivity-domain, def:l-parabolic-neighborhood}
\label{fullmeasure}

Let $({\mathcal M},G)$ be a generalized Ricci flow and let $x\in {\mathcal M}$
with ${\bf t}(x)=T-\tau_1$. Let
 $A\subset {\mathcal
M}\cap {\bf t}^{-1}(-\infty,T-\tau_1)$ be a subset whose intersection with each
time-slice $M_t$ is measurable. Suppose that there is a  subset $F\subset
{\mathcal M}$ such that $|\nabla R|$ and $|\Ric|$ are bounded on $F$ and
such that every minimizing ${\mathcal L}$ geodesic from $x$ to any point in a
neighborhood, $\nu(A)$, of $A$ is contained in $F$. Then for each $\tau\in
(\tau_1,\bar\tau]$ the intersection of $A$ with ${\mathcal U}_x(\tau)$ is an
open subset of full measure in $A\cap M_{T-\tau}$.
\end{corollary}

\begin{proof}
  \uses{cor:l-geodesic-speed-bound, def:l-exponential-map, def:l-injectivity-domain}

\sketch
Since ${\mathcal U}_x(\tau)$ is an open subset of $M_{T-\tau}$, the complement
of $\nu(A)\cap{\mathcal U}_x(\tau)$ in $\nu(A)\cap M_{T-\tau}$ is a closed
subset of $\nu(A)\cap M_{T-\tau}$. Since there is a minimizing ${\mathcal
L}$-geodesic to every point of $\nu(A)\cap M_{T-\tau}$, the ${\mathcal
L}$-exponential map ${\mathcal L}\exp^\tau_x$ is onto $\nu(A)\cap
M_{T-\tau}$.



 According to the next claim, the first subset given in Claim~\ref{claim:complement-injectivity-set-two-types} is contained in the
set of points of $\nu(A)\cap M_{T-\tau}$ where $L_x^\tau$ is
non-differentiable. Since $L_x^\tau$ is a locally Lipschitz function on
$\nu(A)$, this subset is of measure zero in $\nu(A)$; see Rademacher's Theorem
on p. 81 of \cite{EvansGariepy}. The second set is of measure zero by Sard's
theorem. This proves, modulo the next claim, that ${\mathcal U}_x(\tau)\cap A$
is full measure in $A\cap M_{T-\tau}$.




This completes the proof of the Corollary~\ref{cor:reduced-length-full-measure}.
\end{proof}

\begin{claim}[Complement of the injectivity set falls into two types]\label{claim:complement-injectivity-set-two-types}
\uses{def:l-exponential-map, def:l-injectivity-domain}
\label{twosets}

The complement of $\nu(A)\cap {\mathcal U}_{x}(\tau)$
in $\nu(A)$ is contained in the union of two sets: The first is the set of
points $z$ where there is more than one minimizing ${\mathcal L}$-geodesic from
$x$ ending at $z$ and if $Z$ is the initial condition for any minimizing
${\mathcal L}$-geodesic to $z$ then the differential of ${\mathcal L}\mathrm{exp}_x^\tau$ at any $Z$ is an isomorphism. The second is the intersection of
the set of critical values of ${\mathcal L}\exp^\tau$ with $\nu(A)\cap
M_{T-\tau}$.
\end{claim}

\begin{proof}
\uses{cor:l-geodesic-speed-bound}
\sketch
Suppose that $q\in \nu(A)\cap M_{T-\tau}$ is not contained in ${\mathcal U}_x$.
Let $\gamma_Z$ be a minimal ${\mathcal L}$-geodesic  from $x$ to $q$. If the
differential of ${\mathcal L}\exp_x$ is not an isomorphism at $Z$, then
$q$ is contained in the second set given in the claim. Thus, we can assume that
the differential of ${\mathcal L}\exp_x$ at $Z$, and hence ${\mathcal
L}\exp_x$ identifies a neighborhood $\widetilde V$ of $Z$ in ${\mathcal
H}T_z{\mathcal M}$ with a neighborhood $V\subset \nu(A)$  of $q$ in
$M_{T-\tau}$. Suppose that there is no neighborhood $\widetilde V'\subset
\widetilde V$ of $Z$ so that the ${\mathcal L}$-geodesics are unique minimal
${\mathcal L}$-geodesics to their endpoints in $M_{T_\tau}$. Then there is a
sequence of minimizing ${\mathcal L}$-geodesics $\gamma_n$ whose endpoints
converge to $q$, but so that no $\gamma_n$ has initial condition contained in
$\widetilde V'$. By hypothesis all of these geodesics are contained in $F$ and
hence $|\Ric|$ and $|\nabla R|$ are uniformly bounded on these geodesics.
Also,  by the continuity of ${\mathcal L}$, the ${\mathcal L}$-lengths of
$\gamma_n$ are uniformly bounded as $n$ tends to infinity. By
Corollary~\ref{cor:l-geodesic-speed-bound} we see that the initial conditions
$Z_n=\sqrt{\tau_1}X_{\gamma_n}(\tau_1)$ (meaning the limit as $\tau\rightarrow
0$ of these quantities in the case when $\tau_1=0$) are of uniformly bounded
norm. Hence, passing to a subsequence we can arrange that the $Z_n$ converge to
some $Z_\infty$. The tangent vector $Z_\infty$ is the initial condition of an
${\mathcal L}$-geodesic $\gamma_\infty$. Since the $\gamma_n$ are minimizing
${\mathcal L}$-geodesics to a sequence of points $q_n$ converging to $q$, by
continuity it follows that $\gamma_\infty$ is a minimizing ${\mathcal
L}$-geodesic to $q$. Since none of the $Z_n$ is contained in $\widetilde V'$,
it follows that $Z_\infty\not= Z$. This is a contradiction, showing that
throughout some neighborhood $\widetilde V'$ of $Z$ the ${\mathcal
L}$-geodesics are unique minimizing ${\mathcal L}$-geodesics and completing the
proof of the claim.
\end{proof}

\begin{claim}[Two distinct minimizing geodesics give non-differentiability]\label{claim:l-length-nondifferentiable-at-multiple-geodesics}
\uses{lem:l-length-tilde-gradient}

Let $z\in M_{T-\tau}$. Suppose that there is a neighborhood of $z$ in
$M_{T-\tau}$ with the property that every point of the neighborhood is the
endpoint of a minimizing ${\mathcal L}$-geodesic from $x$, so that $L^\tau_x$
is defined on this neighborhood of $z$. Suppose that there are two distinct,
minimizing ${\mathcal L}$-geodesics $\gamma_{Z_1}$ and $\gamma_{Z_2}$ from $x$
ending at $z$ with the property that the differential of ${\mathcal L}\mathrm{exp}^{\tau}$ is an isomorphism at both $Z_1$ and $Z_2$. Then the function
$L^\tau_x$ is non-differentiable at $z$.
\end{claim}

\begin{proof}
\sketch
Suppose that
$\gamma_{Z_0}|_{[0,\tau]}$ is an ${\mathcal L}$-minimal ${\mathcal L}$-geodesic
and that $d_{Z_0}{\mathcal L}\exp_x^\tau$ is an isomorphism. Then use
${\mathcal L}\exp_x^{\tau}$ to identify a neighborhood of $Z_0\in T_xM$
with a neighborhood of $z$ in $M_{T-\tau}$, and push  the function
$\widetilde{\mathcal L}^\tau_x$ on this neighborhood of $Z_0$ down to a
function $L_{Z_0}$ on a neighborhood in $M_{T-\tau}$ of $z$. According to
Lemma~\ref{lem:l-length-tilde-gradient} the resulting function $L_{Z_0}$ is smooth and its gradient
at $z$ is $2\sqrt{\tau}X(\tau)$. Now suppose that there is a second ${\mathcal
L}$-minimizing ${\mathcal L}$-geodesic to $z$ with initial condition
$Z_1\not=Z_0$ and with $d_{Z_1}{\mathcal L}\exp_x^\tau$ being an
isomorphism. Then near $z$ the function $L_x^{\tau}$ is less than or equal to
the minimum of two smooth functions $L_{Z_0}$ and $L_{Z_1}$. We have
$L_{Z_0}(z)=L_{Z_1}(z)=L_x^\tau(z)$,  and furthermore, $L_{Z_0}$ and $L_{Z_1}$
have distinct gradients at $z$. It follows that $L_x^{\tau}$ is not smooth at
$z$.
\end{proof}


\section{Reduced volume}

Here, we  assume that $x\in M_T\subset {\mathcal M}$, so that
$\tau_1=0$ in this subsection.


\begin{definition}[Reduced volume]\label{def:reduced-volume}
\uses{def:reduced-length, def:l-injectivity-domain}
\label{redvol}

Let $A\subset {\mathcal U}_x(\tau)$ be a  measurable subset of
$M_{T-\tau}$. The $\mathcal{L}$-{\em reduced volume of $A$ from $x$}
(or the {\em reduced volume} for short) is
defined to be
$$\widetilde V_x(A) = \int_{A} \tau^{-\frac{n}{2}}\exp(-l_x(q))dq$$ where $dq$ is the volume
element of the metric $G(T-\tau)$.
\end{definition}




\begin{lemma}[Reduced volume as a pullback integral]\label{lem:reduced-volume-pullback}
\uses{def:reduced-volume, def:l-exponential-map, def:reduced-length-tangential}
\label{redvoltilde}

Let $A\subset {\mathcal U}_x(\tau)$ be a  measurable subset. Define
$\widetilde A\subset \widetilde{\mathcal U}_x(\tau)$ to be the
pre-image under ${\mathcal L}\exp_x^{\tau}$ of $A$. Then
$$\widetilde V_x(A)=\int_{\widetilde A}\tau^{-\frac{n}{2}}\exp(-\tilde l(Z, \tau)){\mathcal J}(Z,\tau)dZ,$$ where $dZ$ is the usual
Euclidean volume element and ${\mathcal J}(Z,\tau)$ is the Jacobian determinant
of ${\mathcal L}\exp_x^{\tau}$ at $Z\in T_xM_T$.
\end{lemma}

\begin{proof}
\sketch
This is simply the change of variables formula for integration.
\end{proof}






Before we can study the reduced volume we must study the function
that appears as the integrand in its definition. To understand the
limit as $\tau\rightarrow 0$ requires a rescaling argument.


\subsection{Rescaling}


Fix $Q>0$. We rescale to form $(Q{\mathcal M},QG)$ and then we shift
the time by $T-QT$ so that the time-slice $M_T$ in the original flow
is the $T$ time-slice of the new flow. We call the result
$({\mathcal M}',G')$. Recall that $\tau=T-{\bf t}$ is the parameter
for ${\mathcal L}$-geodesics in $({\mathcal M},G)$ The corresponding
parameter in the rescaled flow $({\mathcal M}',G')$ is $\tau'=T-{\bf
t'}=Q\tau$. We denote by ${\mathcal L}'\exp_x$ the ${\mathcal
L}$-exponential map from $x$ in $({\mathcal M}',G')$, and by $l_x'$
the reduced length function for this Ricci flow. The associated
function on the tangent space is denoted $\tilde l'$.

\begin{lemma}[${\mathcal L}$-length rescaling]\label{lem:l-length-rescaling}
\uses{def:l-length, def:l-geodesic}
\label{Lrescale}

Let $({\mathcal M},G)$ be a generalized Ricci flow and let $x\in
M_T\subset {\mathcal M}$. Fix $Q>0$ and let $({\mathcal M}',G')$ be
the $Q$ scaling and shifting of $({\mathcal M},G)$ as described in
the previous paragraph. Let $\iota\colon{\mathcal M}\to {\mathcal
M}'$ be the identity map. Suppose that $\gamma\colon [0,\bar\tau]\to
{\mathcal M}$ is a path parameterized by backward time with
$\gamma(0)=x$. Let $\beta\colon [0,Q\bar\tau]\to Q{\mathcal M}$ be
defined by
$$\beta(\tau')=\iota(\gamma(\tau'/Q)).$$
Then $\beta(0)=x$ and $\beta$ is parameterized by backward time in
$({\mathcal M}',G')$, and ${\mathcal L}(\beta)=\sqrt{Q}{\mathcal
L}(\gamma)$. Furthermore, $\beta$ is an ${\mathcal L}$-geodesic if
and only if $\gamma$ is. In this case, if $Z=\lim_{\tau\rightarrow 0}\sqrt{\tau}X_\gamma(\tau)$ then
$\sqrt{Q^{-1}}Z=\lim_{\tau'\rightarrow
0}\sqrt{\tau'}X_\beta(\tau')$
\end{lemma}

\begin{remark}[Norms are preserved under the rescaling identification]\label{rem:rescaling-norm-identity}
\uses{lem:l-length-rescaling}

Notice that $|Z|_G^2=|\sqrt{Q^{-1}}Z|^2_{G'}$.
\end{remark}

\begin{proof}
\sketch
It is clear that $\beta(0)=x$ and that $\beta$ is parameterized by
backward time in $({\mathcal M}',G')$. Because of the scaling of
space and time by $Q$, we have $R_{G'}=R_G/Q$ and
$X_\beta(\tau')=d\iota(X_\gamma(\tau))/Q$, and hence
$|X_\beta(\tau')|^2_{G'}=\frac{1}{Q}|X_\gamma(\tau)|^2_{G}$. A
direct change of variables in the integral then shows that
$${\mathcal L}(\beta)=\sqrt{Q}{\mathcal L}(\gamma).$$
It follows that $\beta$ is an ${\mathcal L}$-geodesic if and only if
$\gamma$ is. The last statement follows directly.
\end{proof}

Immediately from the definitions we see the following:

\begin{corollary}[${\mathcal L}\exp$ and $\widetilde l$ under rescaling]\label{cor:l-exponential-rescaling-identity}
\uses{lem:l-length-rescaling, def:l-exponential-map, def:reduced-length-tangential}
\label{COR}

With notation as above, and with the substitution $\tau'=Q\tau$, for
any $Z\in {\mathcal H}T_x{\mathcal M}$ and any $\tau>0$ we have
$${\mathcal
L}'\exp_x(\sqrt{Q^{-1}}Z,\tau')=\iota({\mathcal L}\mathrm{exp}_x(Z,\tau))$$ and
$$ \tilde l'(\sqrt{Q^{-1}}Z,\tau')=\tilde l(Z,\tau),$$
whenever these are defined.
\end{corollary}




\subsection{The integrand in the reduced volume integral}


Now we turn our attention to the integrand (over $\widetilde {\mathcal U}_x(\tau)$) in the
reduced volume integral. Namely, set
$$f(\tau)=\tau^{-n/2}e^{-\tilde l(Z,\tau)}{\mathcal J}(Z,\tau),$$
where ${\mathcal J}(Z,\tau)$ is the Jacobian
determinant of ${\mathcal L}\exp_x^\tau$ at the point $Z\in \widetilde
U_x(\tau)\subset T_xM_T$. We wish to see that this quantity is invariant under
the rescaling.

\begin{lemma}[Scale invariance of the reduced-volume integrand]\label{lem:reduced-volume-integrand-scale-invariance}
\uses{def:reduced-volume, cor:l-exponential-rescaling-identity}
\label{Qformula}

With the notation as above let ${\mathcal J}'(Z,\tau')$ denote the
Jacobian determinant of  ${\mathcal L}'\exp_x$. Then, with the
substitution $\tau'=Q\tau$, we have
$$(\tau')^{-n/2}e^{-\tilde l'(\sqrt{Q^{-1}}Z,\tau')}{\mathcal J}'(\sqrt{Q^{-1}}Z,\tau')=
\tau^{-n/2}e^{-\tilde l(Z,\tau)}{\mathcal J}(Z,\tau).$$
\end{lemma}


\begin{proof}
\uses{cor:l-exponential-rescaling-identity,lem:l-length-rescaling}
\sketch
It follows from the first equation in Corollary~\ref{cor:l-exponential-rescaling-identity} that
$$J(\iota){\mathcal J}(Z,\tau)=J(\sqrt{Q^{-1}}){\mathcal
J}'(\sqrt{Q^{-1}}Z,\tau'),$$ where $J(\iota)$ is the Jacobian
determinant of $\iota$ at ${\mathcal L}\exp_x(Z,\tau)$ and
$J(\sqrt{Q^{-1}})$ is the Jacobian determinant of multiplication by
$\sqrt{Q^{-1}}$ as a map from $T_xM_T$ to itself, where the domain
has the metric $G$ and the range has metric $G'=QG$. Clearly, with
these conventions, we have $J(\iota)=Q^{n/2}$ and
$J(\sqrt{Q^{-1}})=1$. Hence, we conclude
$$Q^{n/2}{\mathcal J}(Z,\tau)={\mathcal J}'(\sqrt{Q^{-1}}Z,\tau').$$

Letting $\gamma$ be the ${\mathcal L}$-geodesic in $({\mathcal
M},G)$ with initial condition $Z$ and $\beta$ the ${\mathcal
L}$-geodesic in $({\mathcal M}',G')$ with initial condition
$\sqrt{Q^{-1}}Z$,  by Lemma~\ref{lem:l-length-rescaling} we have
$\gamma(\tau)=\beta(\tau')$. From Corollary~\ref{cor:l-exponential-rescaling-identity} and the
definition of the reduced length, we get $$\tilde
l'(\sqrt{Q^{-1}}Z,\tau')=\tilde l(\gamma,\tau).$$ Plugging these in
gives the result.
\end{proof}

Let us evaluate $f(\tau)$ in the case of $\Ar^n$ with the Ricci flow being the
constant family of Euclidean metrics.

\begin{example}[${\mathcal L}\exp$ in Euclidean space]\label{ex:l-exponential-euclidean}
\uses{def:l-exponential-map, def:reduced-length-tangential}

Let the Ricci flow be the constant family of standard metrics on
$\Ar^n$. Fix $x=(p,T)\in \Ar ^n\times (-\infty,\infty)$. Then
$${\mathcal L}\exp_x(Z,\tau)=(p+2\sqrt{\tau}Z,T-\tau).$$
In particular, the Jacobian determinant of ${\mathcal L}\mathrm{exp}_{(x,T)}^\tau$ is constant and equal to $2^n\tau^{n/2}$. The
$\widetilde l$-length of the ${\mathcal L}$-geodesic
$\gamma_Z(\tau)=(p+2\sqrt{\tau}Z, T-\tau),\ 0\le \tau\le \bar\tau$,
is $|Z|^2$.
\end{example}

Putting these computations together gives the following.

\begin{claim}[The reduced-volume integrand in Euclidean space]\label{claim:reduced-volume-integrand-euclidean}
\uses{ex:l-exponential-euclidean}
\label{Arnquantity}

In the case of the constant flow on Euclidean space we have
$$f(\tau)=\tau^{-n/2}e^{-\tilde l(Z,\tau)}{\mathcal J}(Z,\tau)=2^ne^{-|Z|^2}.$$
\end{claim}

This computation has consequences for all Ricci flows.

\begin{proposition}[Limit of the reduced-volume integrand]\label{prop:reduced-volume-integrand-limit}
\uses{claim:reduced-volume-integrand-euclidean, def:l-exponential-map}
\label{Jlimit}

Let $({\mathcal M},G)$ be a generalized Ricci flow and let $x\in M_T\subset
{\mathcal M}$. Then, for any $A<\infty$, there is $\delta>0$ such that the map
${\mathcal L}\exp_x$ is defined on $B(0,A)\times (0,\delta)$, where
$B(0,A)$ is the ball of radius $A$ centered at the origin in $T_xM_T$.
Moveover, ${\mathcal L}\exp_x$ defines a diffeomorphism of $B(0,A)\times
(0,\delta)$ onto an open subset of ${\mathcal M}$. Furthermore,
$$\lim_{\tau\rightarrow 0}\tau^{-n/2}e^{-\tilde l(Z,\tau)}{\mathcal J}(Z,\tau)=2^ne^{-|Z|^2},$$
where the convergence is uniform on each compact subset of $T_xM_T$.
\end{proposition}

\begin{proof}
\uses{lem:reduced-volume-integrand-scale-invariance,claim:reduced-volume-integrand-euclidean,lem:l-length-rescaling}
\sketch
First notice that since $T$ is greater than the initial time of ${\mathcal M}$,
there is $\epsilon>0$, and an embedding $\rho\colon B(x,T,\epsilon)\times
[T-\epsilon,T]\to {\mathcal M}$ compatible with time  and the vector field. By
taking $\epsilon>0$ smaller if necessary, we can assume that the image of
$\rho$ has compact closure in ${\mathcal M}$. By compactness every higher
partial derivative (both spatial and temporal) of the metric is bounded on the
image of $\rho$.

Now take a sequence of positive constants $\tau_k$ tending to $0$ as
$k\rightarrow\infty$, and set $Q_k=\tau_k^{-1}$. We let $({\mathcal
M}_k,G_k)$ be the $Q_k$-rescaling and shifting of $({\mathcal M},G)$
as described at the beginning of this section. The rescaled version
of $\rho$ is an embedding
$$\rho_k\colon B_{G_k}(x,T,\sqrt{Q_k}\epsilon)\times [T-Q_k\epsilon,T]\to
{\mathcal M}_k$$ compatible with the time function ${\bf t}_k$ and the vector
field. Furthermore, uniformly on the image of $\rho_k$, every higher partial
derivative of the metric is bounded by a constant that goes to zero with $k$.
Thus, the generalized Ricci flows $({\mathcal M}_k,G_k)$ based at $x$ converge
geometrically to the constant family of Euclidean metrics on $\Ar^n$. Since the
ODE given in Equation~(\ref{reparam}) is regular even at $0$, this implies that
the ${\mathcal L}$-exponential maps for these flows converge uniformly on the
balls of finite radius centered at the origin of the tangent spaces at $x$ to
the ${\mathcal L}$-exponential map of $\Ar^n$ at the origin. Of course, if
$Z\in T_xM_T$ is an initial condition for an ${\mathcal L}$-geodesic in
$({\mathcal M},G)$, then $\sqrt{Q_k^{-1}}Z$ is the initial condition for the
corresponding ${\mathcal L}$-geodesic in $({\mathcal M}_k,G_k)$. But
$|Z|_G=|\sqrt{Q_k^{-1}}Z|_{G_k}$, so that if $Z\in B_G(0,A)$ then
$\sqrt{Q_k^{-1}}Z\in B_{G_k}(0,A)$.
 In particular, we see
that for any $A<\infty$, for all $k$ sufficiently large, the
${\mathcal L}$-geodesics are defined on $B_{G_k}(0,A)\times (0,1]$
and the image is contained in the image of $\rho_k$. Rescaling shows
that for any $A<\infty$  there is $k$ for which the ${\mathcal
L}$-exponential map is defined on $B_G(0,A)\times (0,\tau_k]$ and
has image contained in $\rho$.

Let $Z\in B_Q(0,A)\subset T_xM_T$, and let $\gamma$ be the
${\mathcal L}$-geodesic with $\lim_{\tau\rightarrow
0}\sqrt{\tau}X_\gamma(\tau)=Z$. Let $\gamma_k$ be the corresponding
${\mathcal L}$-geodesic in $({\mathcal M}_k,G_k)$. Then $\lim_{\tau\rightarrow
0}\sqrt{\tau}X_{\gamma_k}(\tau)=\sqrt{\tau_k}Z=Z_k$. Of course,
$|Z_k|^2_{G_k}=|Z|^2_G$, meaning that $Z_k$ is contained in the ball
$B_{G_k}(0,A)\subset T_xM_T$ for all $k$. Hence, by passing to a
subsequence we can assume that, in the geometric limit, the
$\sqrt{\tau_k}Z$ converge to a tangent vector $Z'$ in the ball of
radius $A$ centered at the origin in the tangent space to Euclidean
space. Of course $|Z'|^2=|Z|_G^2$. By Claim~\ref{claim:reduced-volume-integrand-euclidean}, this
means that we have
$$\lim_{k\rightarrow \infty}1^{-n/2}e^{-\tilde l_k(\sqrt{Q_k^{-1}}Z,1)}{\mathcal
J}_k(\sqrt{Q_k^{-1}}Z,1)=2^ne^{-|Z|^2},$$ where ${\mathcal J}_k$ is
the Jacobian determinant of the ${\mathcal L}$-exponential map for
$({\mathcal M}_k,G_k)$.
 Of course, since
$\tau_k=Q_k^{-1}$, by Lemma~\ref{lem:reduced-volume-integrand-scale-invariance} we have
$$1^{-n/2}e^{-\tilde l_k(\sqrt{Q_k^{-1}}Z,1)}{\mathcal
J}_k(\sqrt{Q_k^{-1}}Z,1)=\tau_k^{-n/2}e^{-\tilde
l(Z,\tau_k)}{\mathcal J}(Z,\tau_k).$$ This establishes the limiting
result.

Since the geometric limits are uniform on balls of finite radius
centered at the origin in the tangent space, the above limit also is
uniform over each of these balls.
\end{proof}

\begin{corollary}[The injectivity set contains a fixed ball for small $\tau$]\label{cor:injectivity-set-contains-ball-small-tau}
\uses{prop:reduced-volume-integrand-limit}
\label{Jlimit1}

Let $({\mathcal M},G)$ be a generalized Ricci flow whose sectional
curvatures are bounded. For any $x\in M_T$ and any $R<\infty$ for
all $\tau>0$ sufficiently small, the ball of radius $R$ centered at
the origin in $T_xM_T$ is contained in $\widetilde {\mathcal
U}_x(\tau)$.
\end{corollary}



\begin{proof}
\uses{prop:reduced-volume-integrand-limit,lem:l-geodesic-min-max-estimate,def:l-exponential-map}
\sketch
According to the last result, given $R<\infty$, for all $\delta>0$ sufficiently
small the ball of radius $R$ centered at the origin in $T_xM_T$ is contained in
${\mathcal D}_x^\delta$, in the domain of definition of ${\mathcal L}\mathrm{exp}_x^\delta$ as given in Definition~\ref{def:l-exponential-map}, and ${\mathcal L}\mathrm{exp}_x$ is a diffeomorphism on this subset. We shall show that if $\delta>0$ is
sufficiently small, then the resulting ${\mathcal L}$-geodesic $\gamma$ is the
unique minimizing ${\mathcal L}$-geodesic. If not then there must be another,
distinct ${\mathcal L}$-geodesic to this point whose ${\mathcal L}$-length is
no greater than that of $\gamma$. According to Lemma~\ref{lem:l-geodesic-min-max-estimate} there is a
constant $C_1$ depending on the curvature bound and on $\delta$ such that if
$Z$ is an initial condition for an ${\mathcal L}$-geodesic then for all
$\tau\in (0,\delta)$ we have
$$C_1^{-1}\left(|Z|-\frac{(C_1-1)}{2}\sqrt{\delta}\right)\le \sqrt{\tau}|X(\tau)|\le
C_1|Z|+\frac{(C_1-1)}{2}\sqrt{\delta}.$$ From the formula given in
Lemma~\ref{lem:l-geodesic-min-max-estimate} for $C_1$, it follows that, fixing the bound of
the curvature and its derivatives,
 $C_1\rightarrow 1$ as $\delta\rightarrow 0$. Thus, with
a given curvature bound, for $\delta$ sufficiently small,
$\sqrt{\tau}|X(\tau)|$ is almost a constant along ${\mathcal
L}$-geodesics. Hence, the integral of $\sqrt{\tau}|X(\tau)|^2$ is
approximately $2\sqrt{\delta}|Z|^2$. On the other hand, the absolute
value of the integral of $\sqrt{\tau}R(\gamma(\tau))$ is at most
$2C_0\delta^{3/2}/3$ where $C_0$ is an upper bound for the absolute
value of the scalar curvature.

Given $R<\infty$, choose  $\delta>0$ sufficiently small such that ${\mathcal
L}\exp_x$ is a diffeomorphism on the ball of radius $9R$ centered at the
origin and such that the following estimate holds: The ${\mathcal L}$-length of
an ${\mathcal L}$-geodesic defined on $[0,\delta]$ with initial condition $Z$
is between $\sqrt{\delta}|Z|^2$ and $3\sqrt{\delta}|Z|^2$. To ensure the latter
estimate we need only take $\delta$ sufficiently small given the curvature
bounds and the dimension. Hence, for these $\delta$ no ${\mathcal L}$-geodesic
with initial condition outside the ball of radius $9R$ centered at the origin
in $T_xM_T$ can be as short as any ${\mathcal L}$-geodesic with initial
condition in the ball of radius $R$ centered at the same point. This means that
the ${\mathcal L}$-geodesics defined on $[0,\delta]$ with initial condition
$|Z|$ with $|Z|<R$ are unique minimizing ${\mathcal L}$-geodesics.
\end{proof}

\subsection{Monotonicity of reduced volume}

Now we are ready to state and prove our main result concerning the
reduced volume.




\begin{theorem}[Monotonicity of reduced volume]\label{thm:reduced-volume-monotone}
\uses{def:reduced-volume}
\label{Amono}

Fix $x\in M_T\subset {\mathcal M}$. Let $A\subset{\mathcal
U}_x\subset {\mathcal M}$ be an open  subset. We suppose that for
any $0<\tau\le \bar\tau$ and any $y\in A_\tau=A\cap M_{T-\tau}$ the
minimizing ${\mathcal L}$-geodesic from $x$ to $y$ contained in
$A\cup \{x\}$.  Then $\widetilde V_x(A_\tau)$ is a non-increasing
function of $\tau$ for all $0<\tau\le \bar\tau$.
\end{theorem}





\begin{proof}
  \uses{lem:reduced-volume-pullback, def:reduced-length-tangential, prop:reduced-volume-integrand-limit, lem:l-length-hessian-jacobi-formula, prop:l-length-hessian-inequality, prop:l-length-laplacian-inequality, cor:reduced-length-tangential-k-derivative, claim:h-sum-formula}

\sketch
Fix $\tau_0\in (0,\bar\tau]$. To prove the theorem we shall show
that for any $0<\tau<\tau_0$ we have $\widetilde V_x(A_\tau)\ge
\widetilde V_x(A_{\tau_0})$. Let $\widetilde A_{\tau_0}\subset
\widetilde {\mathcal U}_x(\tau_0)$ be the pre-image under ${\mathcal
L}\exp_x^{\tau_0}$ of $A_{\tau_0}$. For each $0<\tau\le \tau_0$
we set
$$A_{\tau,\tau_0}={\mathcal L}\exp_x^\tau(\widetilde
A_{\tau_0})\subset M_{T-\tau}.$$ It follows from the assumption on
$A$ that $A_{\tau,\tau_0}\subset A_\tau$, so that $\widetilde
V_x(A_{\tau,\tau_0})\le \widetilde V_x(A_\tau)$. Thus, it suffices
to show that for all $0<\tau\le \tau_0$ we have
$$\widetilde V_x(A_{\tau,\tau_0})\ge \widetilde V_x(\tau_0).$$

Since
$$ \widetilde V_x(A_{\tau,\tau_0})=\int_{\widetilde A_{\tau_0}} \tau^{-\frac{n}{2}}
\exp(-\tilde l(Z,\tau)){\mathcal J}(Z,\tau)dZ,$$ the theorem
follows from:





As we have already seen, this proposition implies
 Theorem~\ref{thm:reduced-volume-monotone}, and hence the proof of this theorem is complete.
\end{proof}

\begin{proposition}[Pointwise monotonicity of the reduced-volume integrand]\label{prop:reduced-volume-integrand-monotone}
\uses{def:reduced-length-tangential, prop:reduced-volume-integrand-limit}
\label{fmonotone}

For each $Z\in \widetilde
{\mathcal U}_x(\bar\tau)\subset T_xM_T$ the function
$$f(Z,\tau)=\tau^{-\frac{n}{2}}e^{-\tilde l(Z,\tau)}{\mathcal J}(Z,\tau)$$
is a non-increasing function of $\tau$ on the interval
$(0,\bar\tau]$ with $\lim_{\tau\rightarrow
0}f(Z,\tau)=2^ne^{-|Z|^2}$, the limit being uniform on any compact
subset of $T_xM_T$.
\end{proposition}




\begin{proof}
\uses{lem:l-length-hessian-jacobi-formula,prop:l-length-hessian-inequality,prop:l-length-laplacian-inequality,cor:reduced-length-tangential-k-derivative,prop:reduced-volume-integrand-limit,claim:h-sum-formula}
\sketch
 First, we analyze the Jacobian
${\mathcal J}(Z,\tau)$. We know that $\mathcal{L}\exp_x^{\tau}$
is smooth in a neighborhood of $Z$. Choose a basis
$\{\partial_\alpha\}$ for $T_xM_T$ such that $\partial_\alpha$
pushes forward under the differential at $Z$ of ${\mathcal L}\mathrm{exp}_x^\tau$  to an orthonormal basis $\{Y_\alpha\}$ for
$M_{T-\tau}$ at $\gamma_Z(\tau)$. Notice that, letting $\tau'$ range
from $0$ to $\tau$ and taking the push-forward of the
$\partial_\alpha$ under the differential at $Z$ of ${\mathcal L}\mathrm{exp}_x^{\tau'}$ produces a basis of ${\mathcal L}$-Jacobi fields
$\{Y_\alpha(\tau')\}$ along $\gamma_Z$. With this understood, we
have:
\begin{eqnarray*}
\frac{\partial}{\partial \tau}\ln\mathcal{J}|_\tau & = &
\frac{d}{d\tau}\ln(\sqrt{\det(\langle
Y_{\alpha},Y_{\beta}\rangle
)})\\
 & = &
\frac{1}{2}\Bigl(\frac{d}{d\tau}\sum_{\alpha}\abs{Y_{\alpha}}^{2}\Bigr)\Bigl|_{\tau}\Bigr.
.
\end{eqnarray*}
By Lemma~\ref{lem:l-length-hessian-jacobi-formula} and by Proposition~\ref{prop:l-length-hessian-inequality} (recall that
$\tau_1=0$) we have
\begin{eqnarray}\nonumber
\frac{1}{2}\frac{d}{d\tau}|Y_\alpha(\tau)|^2 & = &
\frac{1}{2\sqrt{\tau}}\Hess(L)(Y_\alpha,Y_\alpha) +\Ric(Y_\alpha,Y_\alpha) \\ & \le &
\frac{1}{2\tau}-\frac{1}{2\sqrt{\tau}}\int_0^{\tau}\sqrt{\tau'}H(X,\tilde
Y_\alpha(\tau'))d\tau',\label{basic}\end{eqnarray} where $\tilde
Y_\alpha(\tau')$ is the adapted vector field along $\gamma$ with
$\tilde Y(\tau)=Y_\alpha(\tau)$. Summing over $\alpha$ as in the
proof of Proposition~\ref{prop:l-length-laplacian-inequality} and Claim~\ref{claim:h-sum-formula} yields
\begin{eqnarray}
\frac{\partial}{\partial\tau}\ln\mathcal{J}(Z,\tau)|_{\tau} &
\le & \frac{n}{2\tau} -
\frac{1}{2\sqrt{\tau}}\sum_{\alpha}\int_{0}^{\tau}\sqrt{\tau'}H(X,\tilde
Y_{\alpha(\tau')})d\tau' \nonumber
\\&=  & \frac{n}{2\tau} - \frac{1}{2}\tau^{-\frac{3}{2}}K^\tau(\gamma_Z).\label{Jinequ}
\end{eqnarray}

On $\widetilde {\mathcal U}_x(\tau)$ the expression
$\tau^{-\frac{n}{2}}e^{-\widetilde l(Z,\tau)}{\mathcal J}(Z,\tau)$
is positive, and so we have
$$\frac{\partial}{\partial\tau}\ln\left(\tau^{-\frac{n}{2}}e^{-\widetilde
l(Z,\tau)}{\mathcal J}(Z,\tau)\right)\le
\left(-\frac{n}{2\tau}-\frac{d\widetilde l}{d
\tau}+\frac{n}{2\tau}-\frac{1}{2}\tau^{-\frac{3}{2}}K^\tau(\gamma_Z)\right).$$


Corollary~\ref{cor:reduced-length-tangential-k-derivative} says that the right-hand side of the
previous inequality is zero. Hence, we conclude
\begin{equation}\label{Jineq}
\frac{d}{d\tau}\left(\tau^{-\frac{n}{2}}e^{-\tilde
l(X,\tau)}{\mathcal J}(X,\tau)\right)\le 0. \end{equation}

This proves the inequality given in the statement of the
proposition. The limit statement as $\tau\rightarrow 0$ is contained
in Proposition~\ref{prop:reduced-volume-integrand-limit}.
\end{proof}

Notice that we have established the following:


\begin{corollary}[Reduced volume is bounded]\label{cor:reduced-volume-bounded}
\uses{thm:reduced-volume-monotone}
\label{finiteredvol}

For any measurable subset $A\subset {\mathcal U}_x(\tau)$ the
reduced volume  $\widetilde V_x(A)$ is at most $(4\pi)^{n/2}$.
\end{corollary}


\begin{proof}
\uses{thm:reduced-volume-monotone}
\sketch
Let $\widetilde A\subset \widetilde {\mathcal U}_x(\tau)$ be the pre-image
of $A$. We have seen that
$$\widetilde V_x(A)=\int_{A}\tau^{-n/2}e^{-l(q,\tau)}dq=\int_{\widetilde
A}\tau^{-n/2}e^{-\tilde l(Z,\tau)}{\mathcal J}(Z,\tau)dz.$$ By
Theorem~\ref{thm:reduced-volume-monotone} we see that $\tau^{-n/2}e^{-\tilde
l(Z,\tau)}{\mathcal J}(Z,\tau)$ is a non-increasing function of
$\tau$ whose limit as $\tau\rightarrow 0$ is the restriction of
$2^ne^{-|Z|^2}$ to $\widetilde A$. The result is immediate from
Lebesgue dominated convergence.
\end{proof}
